% \documentclass[final,5p,times,twocolumn,authoryear]{elsarticle}
\documentclass[]{article}
\usepackage[margin=0.6in]{geometry}
\usepackage{amssymb}
\usepackage{amsmath}
\usepackage{lipsum}
\usepackage{booktabs}
\usepackage{caption}
\usepackage{xcolor}
\usepackage{hyperref}
\usepackage{array} % for better column formatting
\usepackage{multirow} % for multirow cells in tables
\usepackage{booktabs}    % for better-looking tables
\usepackage{changepage}  % for adjusting page layout locally
\usepackage{float}
% \usepackage{xeCJK} % For Chinese
\usepackage{url}
\usepackage{tikz}
\usepackage{natbib}
\usepackage{graphicx}
\usepackage{placeins}
\usepackage{float}
\usepackage{indentfirst}
\usepackage{setspace}
\usepackage{rotating}
\usepackage{tabularx}
\usepackage{caption}  % for caption customization
\usepackage{listings} % for listing figures
\usepackage{ragged2e} % for better text alignment
\usepackage{xltabular} % for 'xltabular' environment (similar to 'tabularx' but allows page breaks)
\usepackage{xurl}

% Define new column types with different alignments and widths
\newcolumntype{L}{>{\RaggedRight\arraybackslash}X} % Left-aligned X column
\newcolumntype{C}{>{\Centering\arraybackslash}X}   % Center-aligned X column
\newcolumntype{R}{>{\RaggedLeft\arraybackslash}X}  % Right-aligned X column

\captionsetup[table]{justification=raggedright, singlelinecheck=false}
\hypersetup{colorlinks=true, citecolor=teal, linkcolor=teal, urlcolor=teal}

\newcommand{\kms}{km\,s$^{-1}$}
\newcommand{\msun}{$M_\odot}

\usepackage{abstract}
\renewcommand{\absleftindent}{0mm}  % Adjust the left margin
\renewcommand{\absrightindent}{0mm} % Adjust the right margin

% % Custom command to store figure information
% \newcommand{\figsource}[2]{%
%     \caption[#1]{#1 \\ 
%     Source: #2}%
% }

\begin{document}
\sloppy

\title{Develop the design principles for demand-responsive railway station areas}
\author{}
\date{}

\maketitle

\begin{abstract}
\textbf{Aims}: Station areas suffer from overcrowding and emptiness due to the fluctuation caused by events. This paper aims to establish design principles that help design demand-responsive station areas. \textbf{Methods}: It firstly establishes design objectives and spatial dimensions, then collects data from literature and cases, followed by data processing through systematic examination and research-by-design; It evaluates the initial findings according to design objectives and spatial relevance, and reflects the final findings through experts interview. \textbf{Results}: More than twenty design principles emerged, and they can be viewed from various perspectives. These principles can better facilitate design practice, and the perspectives lead to new solutions being found in the future.
\end{abstract}

% \begin{keyword}
% Overcrowding \sep emptiness \sep urban design \sep pattern language
% \end{keyword}

\providecommand{\keywords}[1]{\textbf{Keywords:} #1}
\keywords{Overcrowding, Emptiness, Urban design, Pattern language}

% \end{frontmatter}

\section{Introduction}
\label{Introduction}
In railway stations, overcrowding and emptiness (Fig. \ref{fig:overcrowdingemptiness}), which are two extremes of passenger fluctuation (Section \ref{The nature of overcrowding and emptiness}), are tricky issues causing economic and experiential losses. Overcrowding is unpleasant for users and, more importantly, is a safety concern of stampedes \citep{chinadailyStampedeKillsOne2008, almeidaHumanStampedesUpdated2019}. Emptiness is a waste of space as resources for society, especially in high-density urban areas or cities; It is a waste of maintenance for railway companies; It causes economic loss, as stores or real estate property owners get less commercial revenue due to lack of customer patronage; It decreases users' experience, as spaces being empty makes people feel dangerous, especially at night. All these negative effects necessitate solutions to address overcrowding and emptiness in railway stations.

\label{General relevance to non-station research topics}
Addressing overcrowding and emptiness is relevant for many stations as they are challenged by various types of events. Such as spring festivals in China, Vietnam, and South Korea; and beer festivals, carnivals, and sports events in European cities \citep{cityrFestivalkalenderRotterdamCity}. 
As events are ever-growing in cities \citep{tallieuEuropeanFestivalsAssociation} enriching the city dynamics, it is necessary to build stations that are capable of facilitating this. Also, addressing the after-event emptiness of space and making full use of space as a resource is helpful in promoting sustainability, especially in high-density cities where space is scarce.


\begin{figure}[h]
    \centering
    \includegraphics[width=0.8\linewidth]{Paper3Imgs_pdfs/overcrowdingemptiness.pdf}
    \caption{Overcrowding and emptiness in railway stations areas}


    \label{fig:overcrowdingemptiness}
\end{figure} \FloatBarrier







\subsection{The nature of overcrowding and emptiness}
\label{The nature of overcrowding and emptiness}

\subsubsection{Viewing stations from the pedestrian flow perspective}
\label{Viewing stations from the pedestrian flow perspective}
A basic unit of pedestrian movement can be seen as pedestrians transfer from their originations to their destinations through some spaces/passages/connections (Fig. \ref{fig:flowunit}) \citep{fruinPedestrianPlanningDesign1971}. This basic unit can be repeatedly seen in various forms in the station area (Fig. \ref{fig:flowunit} b, c; Fig. \ref{fig:bottleneck}). Along the passengers’ transfer journeys, sometimes \textbf{bottlenecks} exist (i.e., station components that have a limited capacity to handle the passenger flow, due to various reasons such as security/ticket checks, level changes, narrow passages, a limited number of trains, and disruptions on the railway lines). Station components that are perfectly working during normal times can become bottlenecks during big events (e.g., national holidays and football matches) when there are surging passengers. To accommodate the detained passengers, `redundant' spaces are needed. These redundant spaces can be overcrowded during peak times while empty during non-peak times. Therefore, overcrowding and emptiness, which may seem two opposite phenomena, actually result from the same reason - fluctuated passenger flow due to events.


\begin{figure}[h]
    \centering
    \begin{minipage}{0.45\textwidth}
        \centering
        \includegraphics[width=\columnwidth]{Paper3Imgs_pdfs/flowunit.pdf}
        \caption{A conceptual basic unit of passenger flow}
        \label{fig:flowunit}
    \end{minipage}\hfill
    \begin{minipage}{0.45\textwidth}
        \centering
        \includegraphics[width=\columnwidth]{Paper3Imgs_pdfs/bottleneck.pdf}
        \caption{The mechanism of overcrowding and emptiness in Beijing West railway station area}
        \label{fig:bottleneck}
    \end{minipage}
\end{figure} \FloatBarrier



\subsubsection{Spatial relevance of the problems}
\label{Spatial relevance of the problems}
With professional backgrounds in architecture and urban design, we saw the spatial relevance of overcrowding and emptiness to spatial configurations. The spatial considerations in these two fields usually include the form of space, layout of elements, accessibility, visibility, and so on, on building and district levels. In cases like Beijing West Railway Station (Fig. \ref{fig:beijingwestboth} a, b), overcrowding and emptiness are associated with the poor connection with the neighborhood, the gigantic non-human size of buildings, messy facilities on the square blocking view, and deteriorating space qualities, a lack of visibility from eye level to see the sunken plaza and atriums, traffic separate the use of walkable spaces, and so on. The station is relatively isolated from the neighboring city areas, and spaces within the station areas are separated by bottlenecks, which decrease the peak time capacity and make the station less attractive for non-transport use during non-peak times. 

\begin{figure}[h]
    \centering
    \includegraphics[width=\linewidth]{Paper3Imgs_pdfs/beijingwestboth.pdf}
    \caption{Spatial relevance shown by an example case of Beijing West Railway Station}
    \label{fig:beijingwestboth}
\end{figure} \FloatBarrier

\subsubsection{Temporal and spatial scales of the problem-causing events}
\label{Temporal and spatial scales of the problem-causing events}
Various events on different temporal scales \citep{carmonaPublicPlacesUrban2010} can influence stations (Fig. \ref{fig:temporaldimension}), such as daily and weekly commuting, multi-weekly sports events, yearly holidays, and multi-decennially Olympics. However, daily and weekly commuting is more of a topic for metro stations than railway stations and is commonly investigated in research on ‘TOD (Transit Oriented Development)’ \citep{ibraevaTransitorientedDevelopmentReview2020}. Low-frequency events such as the Olympics are too rare to find patterns \citep{alexanderPatternLanguageTowns1977}. In contrast, multi-weekly to yearly events are peculiar because they result in relatively drastic pedestrian fluctuation, and it is feasible to find patterns. They are of particular concern to railway stations and have barely been studied.

Events on different spatial scales can all impact stations, whether they happen within station buildings, within station areas, or outside station areas. For example, markets are held in the station hall of Zurich Main Station, and concerts are organized in the station hall of Antwerp Central Station. In station areas, there are concerts/sports events in the multiple stadiums near Bijlmer Arena station and conference events in the Jaarbeurs convention center near Utrecht Central Station, which bring surging passengers during event times. Some other events, not necessarily in the station areas, cause peak flow for stations, such as the Chinese national holidays, the Munich beer festival, and the Zandvoort F1 match.
% (Fig. \ref{fig:Spatial scale of events})

\begin{figure}[h]
    \centering
    \includegraphics[width=0.7\linewidth]{Paper3Imgs_pdfs/temporaldimension.pdf}
    \caption{Temporal scales}
    \label{fig:temporaldimension}
\end{figure} \FloatBarrier
% \begin{figure}[h]
%     \centering
%     \includegraphics[width=0.5\linewidth]{Paper3Imgs_pdfs/pxlaceholder.pdf}
%     \caption{Spatial scales of events}
%     \label{fig:Spatial scale of events}
% \end{figure} \FloatBarrier




\subsection{Demand-responsive as an objective}
\label{Demand-responsive as an objective}

The overcrowding and emptiness issues urge station areas to be `demand-responsive.' Since this paper studies the particular issues - overcrowding and emptiness- which are barely addressed by existing spatial design research, there are no corresponding termed solution packages. Therefore, this paper needs to name a term for possible solutions and enrich the meaning of the term. One ideal term this paper found is ‘demand-responsive’ inspired by the word ‘demand-responsive transport,’ which is a practice ideal \citep{enochWhyDemandResponsive2006} used in the transport service field, describing a flexible mode of transportation that adapts to the demands of its user groups \citep{hunkinDemandresponsiveTransportPolicy2018}. This paper transfers the meaning of ‘demand-responsive’ to the station context regarding \textbf{spatial design} perspective, defining it as \textbf{a quality of station areas that can address the use fluctuation caused by events, reduce potential overcrowding and emptiness, make spaces sufficiently be used during non-peak and peak time, and satisfy users’ needs}. (See also Section \ref{Literature}, and \href{https://github.com/1309928130/Paper3_SupplementaryMaterials/blob/main/SupplementaryMaterials.pdf}{Supplementary materials})

% \subsubsection{Demand of who?}
% \label{Demand of who}



\subsection{Spatial solutions as an objective}

This research hypothesizes that spatial solutions (on district and building levels) can contribute to demand-responsive solutions. Firstly, overcrowding and emptiness sometimes result from poor spatial configuration (Section \ref{Spatial relevance of the problems}), spatial relevance of the problems). Therefore, good spatial solutions will make a change. Secondly, it is common for spatial solutions to resolve non-spatial issues or to enhance non-spatial solutions. Indeed, "a pattern of events cannot be separated from the space where it occurs" \citep[p. 73]{alexanderTimelessWayBuilding1979}. Overcrowding and emptiness in stations are often a mixture of spatial and non-spatial issues, and possible solutions can come from spatial design, among other fields such as transportation engineering, management, and planning. For example, a (spatially) well-designed plaza can serve as a basis for (managerially) organizing leisure events, hence reducing emptiness.


\subsubsection{Address both overcrowding and emptiness}
Spatial solutions should be able to address both overcrowding and emptiness and should not ignore one of them. If only overcrowding or emptiness were considered, there could easily be a paradox of ‘larger' or `smaller' spaces for railway station areas. For example, to reduce overcrowding, stations can be designed with `larger' redundant spaces, but then the emptiness issue during non-peak time will be exacerbated; to reduce emptiness, stations can be designed with `smaller' redundant spaces, but they will be more overcrowded during peak time. To avoid the paradox, design solutions should be able to reduce both overcrowding and emptiness. That is to say, solutions should 1) increase transport (i.e., transfer/waiting) capacity without significantly enlarging spaces that are prone to be empty or 2) increase the leisure use of space without significantly reducing transport capacity.

\subsubsection{Spatial scales of concern: the district and building levels (within 250m radius)}
\label{Spatial scales of concern: the district and building levels }

For issues related to railway stations, typically, there are potential solutions across different spatial scales, and spatial solutions at the district level are less studied (Fig. \ref{fig:spatial_scales_station}). 

On regional or city levels, many studies have investigated planning interventions to improve stations' performance \citep{yinEffectsHighspeedRailway2015, bergEuropeanHighspeedTrain1998, borghetti15MinStationCase2021}. In these studies, stations are simplified as nodes within networks. These nodes (station areas) are usually defined with a range \textbf{around 500-1000m} \citep{zhangStudyRevisingSpatioTemporal2019, borghetti15MinStationCase2021}. At these levels, overcrowding can potentially be solved by the shared capacity from other stations \citep{besinovicResilienceRailwayTransport2020, luRailwayVulnerabilityResilience2022} or other mobilities \citep{martinsIndicatorbasedMethodologyAssessing2019}, and emptiness can potentially be solved by better accessibility \citep{borghetti15MinStationCase2021, jehleConnectingPeoplePlaces2022} that promotes more use. 

However, interventions at these higher levels are usually non-spatial, missing the complexity of overcrowding and emptiness relating to 3D and configurations (Section \ref{Spatial relevance of the problems}).

In contrast, district-level interventions can take district structures and 3D urban forms into consideration. And yet, little is known about district-level interventions.

In this research, we investigate station areas with a range within approximately 250-300m, smaller than the typical 500m range in planning-oriented research. This is because we take user perception into consideration. As noted by \citet{peekCreatingSynergyStations2002}, in practice, a 300m radius for the so-called “station environment” is used. This smaller range serves this study's explorative nature to investigate spatial configurations that lie in urban design (and architecture, as explained in the following paragraph) domains. During field visits of selected cases, we found that the stations are usually visually perceivable around the maximum 250m range (see \href{https://github.com/1309928130/Paper3_SupplementaryMaterials/blob/main/SupplementaryMaterials.pdf}{Supplementary Materials}, figs. \ref{fig:station_area_Utrecht}, \ref{fig:station_area_Rotterdam}, and \ref{fig:station_area_BeijingWest}, depicting areas of different stations).  




\begin{figure}[h]
    \centering
    \includegraphics[width=0.6\linewidth]{Paper3Imgs_pdfs/station_area_scales.pdf}
    \caption{Spatial scales of station areas}
    \label{fig:spatial_scales_station}
\end{figure} \FloatBarrier



Besides the district level, the building level (at which we analyze the station building) should also be considered. In railway development, urban design schemes at the district level are commonly determined by the layouts of station buildings at the building level. For example, in Utrecht Central Station and Rotterdam Central Station, the large structures - city axes, are much determined by the qualities of city passages embedded in the station buildings.

Building-level and district-level interventions can be mixed. For example, a book that is highly referenced by railway station studies in China, titled `Holistic Design on Urban Architecture' \citep{hanHolisticDesignUrban1999}, conveys many ideas that mix the boundaries of urban and architectural design, including `the socialization and giantization of architecture,' `the three-dimensionalization and interiorization of urban design.’


Yet Building-level and district-level interventions also have their own strengths and limitations. District-level concepts can be expensive to realize within buildings or lead to non-human scale mega buildings. For example, in Fig. \ref{fig:Archi_urban_proposals_Sloterdijk}, to improve the fragmented Amsterdam Sloterdijk station areas, the architectural proposals featured mega-building roofs that span several hundred meters, while the urban design proposals are more realistic, sorting out the clear human-scale district street structures, with normal-sized buildings.

Therefore, in this research, we investigate the station areas with a range of within 250m at the district and building levels.

\begin{figure}[h]
    \centering
    \includegraphics[width=0.5\linewidth]{Paper3Imgs_pdfs/Sloterdijk_Arch_UD.pdf}
    \caption{Architectural and urban design proposals for Amsterdam Sloterdijk station}
    \label{fig:Archi_urban_proposals_Sloterdijk}
\end{figure} \FloatBarrier

\subsection{Relevance with station-city integration}

Demand-responsive is interconnected with the research topic of `\textbf{station-city integration (SCI)},’ as they can interchangeably be the context or solutions for each other. SCI is a research topic that aims to get various possible benefits by integrating stations and cities. It has a relatively long history of rich academic conversations across multiple fields \citep{chengSymposiumForumChina2021, yinEffectsHighspeedRailway2015}, including spatial design \citep{calthorpeNextAmericanMetropolis1993, dewildeMultipleUseSpace2006, shikataIntegratedStationCityDevelopment2013, zempGenericFunctionsRailway2011, duPlaceQualityHighspeed2021, conceicaoCityStationStation2015}. Spatial interpretations for integration typically include mode integration, mixed use of land, high-density development, the stack of building floors, and so on.

SCI brings potential as well as risks for demand-responsive solutions. When stations and cities are `integrated,' it is possible for stations to benefit from an enlarged accommodation capacity and increased user numbers, therefore relieving the overcrowding and emptiness issues. The opposite side effect is also possible, though; stations may be more overcrowded in peak time due to the increased number of passengers or more empty if the city areas are mostly empty by themselves. The city users may even cause flow conflicts for transport passengers. The city areas can benefit from stations’ massive passengers, a vibrant environment, and increased commercial revenue; it is also possible to have side effects such as being rendered into chaotic areas due to traffic elements. Developing city areas beyond pure stations may require a heavy initial investment, resulting in high maintenance costs after development. The dual possibilities of benefits and side effects make SCI a fruitful and challenging topic for demand-responsive solutions. To reduce overcrowding and emptiness, stations should be integrated with cities carefully, possibly with the help of a systematic assessment of design interventions.



\subsection{Knowledge Gap and Research Question}
\label{Knowledge gap and research question}

The above introduction explains the nature of overcrowding and emptiness (O\&E) as related to events, and reveals the research objective: spatial solutions for demand-responsive (DR) station areas, which bring potential and avoid risks to station-city integration (SCI). 

Aimed at the research objectives, an ongoing Ph.D project is being conducted, which 1) articulates the overcrowding \& emptiness problems from the spatial design perspective, 2) establishes an assessment framework for design evaluation, and 3) develops \textbf{design principles} (fig. \ref{fig:positionP3inPhD}). This paper is part of the project and specifically addresses the design principles. 

The main research question of this paper is: \textbf{What are the design principles for demand-responsive railway station areas?} A subsequent question is:\textbf{ What is a practical way to develop design principles?} 

The remainder of this paper is structured as follows: Chapter \ref{Methodology} describes the methodology, Chapter \ref{Results} reports the results, Chapter \ref{Discussion} discusses the whole research and results, and Chapter \ref{Conclusions} draws conclusions.

\begin{figure}[h]
    \centering
    \includegraphics[width=0.5\columnwidth]{Paper3Imgs_pdfs/positionP3inPhD.pdf}
    \caption{This paper is part of a PhD project at the intersection of three topics}
    \label{fig:positionP3inPhD}
\end{figure} \FloatBarrier

\section{Methodology}
\label{Methodology}
In this section, we first define what `design principles' are, and then describe the methodology for developing them.

\subsection{What is a design principle?}
\label{section:What is a design principle?}

Design principles/guidelines are commonly used in the design field. Their meanings vary in different research or design projects, from vague to more specific \citep{liuMappingLandscapeSpaces2020, vanderhoevenSignificanceUndergroundExperience2016}. We adopt the well-defined concept called `patterns' from \citet{alexanderTimelessWayBuilding1979, alexanderPatternLanguageTowns1977} as what we mean by `design principles.' In the following two paragraphs, we briefly describe the related concept of `patterns' and how we make use of them.

``\textit{Each \textbf{pattern} is a three-part rule, which expresses a relation between a certain \textbf{context, a problem, and a solution}. Each pattern depends both on the smaller patterns it contains, and on the larger patterns within which it is contained}" \citep[pp. 247, 312]{alexanderTimelessWayBuilding1979}. Following these ideas, in this research, we treat the design principles as `solutions' for `problems' within `contexts' (Section \ref{Design principles}). In chapter \ref{Introduction} \hyperref[Introduction]{Introduction}, we already describe the general context and problems. In chapter \ref{Results} \hyperref[Results]{Results}, Section \ref{Design principles} \hyperref[Design principles]{Design principles}, some design principles aimed at their own and more specific problems within certain contexts; we describe these problems and contexts whenever necessary in the description of each principle.

Alexander organizes different patterns in a scalar hierarchical \textbf{network}, in which the relationships between patterns are illustrated (and the network is indeed the structure of the \textbf{pattern language}) \citep[p. 314]{alexanderTimelessWayBuilding1979}. Similarly, we organize all our design principles in a network (fig. \ref{fig:Pattern_language_network}).

\subsection{A seven-step methodology framework for developing design principles}
\label{A seven-step methodology framework for developing design principles}

To systematically discover and evaluate design principles, we propose a methodology seven-step framework (Fig. \ref{fig:methodologyframework}), which includes (S1) developing design objectives and design dimensions, (S2) collecting data from cases and literature, (S3) data processing, (S4) proposing hypotheses, (S5) evaluation of hypotheses, and (S6) presenting final results. Iteration, as a typical procedure in design-related research \citep{nijhuisDesignrelatedResearchLandscape2012}, is also embedded as Step 7 (S7) and as a sub-step within Step 3. The following sections describe some key components of the framework.

\begin{figure*}[htb]
    \centering
    \includegraphics[width=\linewidth]{Paper3Imgs_pdfs/methodologyframework.pdf}
    \caption{The process of developing design principles}
    \label{fig:methodologyframework}
\end{figure*} \FloatBarrier

\subsection{Design objectives}
\label{Design objectives}
Multiple design objectives are established in this paper, which are categorized into a) problem-solving and value-creating and b) user-based. The objectives in the first category are more relevant to stakeholders' interests, and the objectives in the second are more relevant to users’ experiences.

\subsubsection{Problem-solving and value-creating objectives}
\label{Problem-solving and value-creating objectives}
Four design objectives are defined in this category. The initial problems this paper investigates are overcrowding and emptiness, so 1) reducing overcrowding and 2) reducing emptiness naturally becomes the two primary objectives. Since this research is also embedded in the station-city integration topic, it considers beyond the station, setting two extra objectives for the city: 3) adding value to the city and 4) reducing interference to the city. Objectives 1) and 2) are the added values to stations, and 3) and 4) are the added values to cities. In this sense, these objectives contribute to station-city integration. These design objectives serve as guidance for collecting relevant data (in Step 2. Data collection), guide the examination of input data (in Step 3. Data processing), and function as criteria to evaluate the initial results (in Step 5. Evaluation, fig. \ref{fig:methodologyframework}).

\subsubsection{User-based objectives}
\label{User-based objectives}
Various user-based design objectives are set. As transport infrastructure, stations have their primary users as transport passengers; While stations are also promoted as places in the city \citep{bertoliniCitiesRailsRedevelopment1998}, they should also consider the non-passengers (e.g., neighborhood residents) as users. For passengers, according to \citet{vanhagenWaitingExperienceTrain2011}, there are five types of needs: safety, speed, ease, comfort, and experience. A further study by \citet{chenEvaluateUserSatisfaction2024} translated these needs into quantifiable spatial and behavioral/perceptual features. For non-passengers, since they are diverse in nature, it is hard to define their universal needs; One attempt by \citet{carrPublicSpace1992a} specifies the general public space users' needs as comfort, relaxation, passive engagement with the environment, active engagement with the environment, and discovery. This paper therefore set corresponding design objectives as: 5) enhancing the five types of passengers’ needs, and 6) enhancing the multiple general public space needs.

\subsection{Dimensions of spatial and managerial interventions}
\label{Dimensions of spatial and managerial interventions}
This paper examines spatial interventions relevant to design and management. Railway station development is a complex topic addressed by different practice fields, including planning, design, management, and so on. These fields typically have interventions taken on certain temporal scales. Planning strategies may be implemented once a few decades, design and engineering renewals may happen every few years, while management interventions can be taken on a daily basis \citep{carmonaPublicPlacesUrban2010}. As design research, this paper primarily investigates spatial design solutions; while it also investigates solutions relevant to management (and have spatial implications) because events (both as problems and solutions) are closely related to management. Planning is a spatially much broader and temporally long-term field, which is less relevant to the issues this paper addresses and is left out here.

\subsubsection{Spatial scales of station areas}
\label{Spatial scales of station areas}
This paper investigates the station areas at the building and district levels (Fig. \ref{fig:configuration}; Fig. \ref{fig:configurationflow} a). Spatial scale is a critical issue in both railway station research and urban design research \citep{duPlaceQualityHighspeed2021, lloyd-jonesDesignProblems2001}. at the district level, various ranges have been proposed to define railway station areas \citep{calthorpeNextAmericanMetropolis1993, bertoliniStationAreasNodes2008}), typically around 500-700m. This paper, however, chooses a more humble range, 250-300m, what \citet{peekCreatingSynergyStations2002}) call "station environment.” This smaller range is more appropriate as it relates more to user experiential perception \citep{chenEvaluateUserSatisfaction2024}, and is more relevant to the urban design interventions in station development (e.g., see the Rotterdam Central Station redevelopment site range in \href{https://www.archdaily.com/447649/rotterdam-centraal-team-cs}{https://www.archdaily.com}).

\subsubsection{The configuration of station areas}
\label{The configuration of station areas}
This paper views the configuration of station areas from four perspectives to encompass a comprehensive understanding of this complex system. Firstly, in the introduction part (Section \ref{The nature of overcrowding and emptiness}), this paper already explained the station configuration as it facilitates pedestrian flow transfer from originations to destinations (Fig. \ref{fig:flowunit}, \ref{fig:bottleneck}). Secondly, similar to the general flow analysis, designers in design practice often examine specific paths of users to gain experiential/perceptual understandings of space (Fig. \ref{fig:configurationflow} c). Thirdly, design is also commonly treated as the study of elements and layout \citep{dejongeWAYSStudyResearch2002, liuMappingLandscapeSpaces2020}. For station areas at the district level, elements include the station building, squares, city passage, urban furniture, and so on; and the layout is about the composition/constellation/organization/network of these elements. at the building level, the layout is about the organization of various mobility sites (railway, metro, bus, tram, bike, taxi/car) and building components (gates, indoor public spaces, hall, escalators, and so on) (Fig. \ref{fig:configurationflow} b; Fig. \ref{fig:configuration} b). Fourthly, the station areas can also be seen as things that belong to either the station or the neighboring city (Fig. \ref{fig:configurationflow} a). The configuration analyses from these four perspectives can be conducted at the district and the building levels.

\begin{figure*}[h]
    \centering
    \includegraphics[width=0.75\linewidth]{Paper3Imgs_pdfs/configuration.pdf}
    \caption{The station at the district and building levels}
    \label{fig:configuration}
\end{figure*}
\FloatBarrier


\begin{figure*}[h]
    \centering
    \includegraphics[width=\linewidth]{Paper3Imgs_pdfs/configurationflow.pdf}
    \caption{The configuration of the station}
    \label{fig:configurationflow}
\end{figure*}
\FloatBarrier

\subsubsection{Managerial interventions}
\label{Managerial interventions} 
This paper considers two types of managerial interventions: organizing events and regulating access. According to \citet{carrPublicSpace1992a}, there are several types of rights in public space, including access, claim, change, ownership and disposition, and so on. Management interventions can regulate these rights. Inspired by this, the paper realizes at least two types of managerial interventions - organizing events and regulating access (or routes and paths), can help to address overcrowding and emptiness. Firstly, although some (problem-causing) events affecting transport are the sources of the overcrowding and emptiness issues (Section \ref{The nature of overcrowding and emptiness}), some other leisure events can be part of the solutions; When the redundant spaces in station areas are empty during normal times, (solution-oriented) leisure events can be organized to reduce emptiness, such as exhibitions in station buildings, and markets or crowd dancing on station plazas \citep{nieuwenhuizeSalsaDoeGratis2023}. Secondly, regulating access/route/path is commonly seen in crowd management practices to (re-)distribute the pedestrians and avoid overcrowding.

\subsection{Data sources}
\label{Data sources}
This paper uses real-world cases and literature as data sources to discover initial relevant solutions (that already exist in practice). It uses `research by design (RbD) (\ref{fig:researchbydesign})' to further develop new solutions (that do not exist in practice \citep{nijhuisDesignrelatedResearchLandscape2012}). Combining literature and cases with RbD enables us to find solutions from the past and for the future and to leverage both theoretical and practical knowledge (Fig. \ref{fig:datacompare}). 

\begin{figure}[h]
    \centering
    \includegraphics[width=0.3\columnwidth]{Paper3Imgs_pdfs/datacompare.pdf}
    \caption{Data sources}
    \label{fig:datacompare}
\end{figure} \FloatBarrier 

\subsubsection{Cases}
\label{Cases}
The cases used in this paper include seven initial station cases, multiple extended station cases, and non-station cases. Seven initial station cases are selected, including Beijing West Railway Station (RS), Beijing South RS, Shanghai RS, Rotterdam Central Station (CS), Utrecht CS, Bijlmer Arena RS, and Zandvoort RS. They differ in 1) fluctuation contexts and/or 2) performance addressing fluctuation. Some of them - `poor' examples (Beijing West RS and Beijing South RS) - have overcrowding and emptiness phenomena yet not solved, while others - `good' examples - potentially solved these phenomena. Most of them are major stations of the cities (see \href{https://github.com/1309928130/Paper3_SupplementaryMaterials/blob/main/SupplementaryMaterials.pdf}{Supplementary Materials}, Table \ref{tab:case_selected}; Section \ref{Station typology implied by the cases} \hyperref[Station typology implied by the cases]{Station typology implied by the cases}). 
% Hypothetically simulating massive passengers in Rotterdam Central Station, there will be less overcrowding and emptiness.

These initial cases are also chosen due to our physical vicinity to the cases, which enables us to do site visits and get experiential/perceptual heuristics. 

Extended station cases and non-station cases are also added in several situations: 1) When initial station cases are not enough to illustrate the discovered solutions, 2) When extra cases extend the generalizability of the design principles. The non-station cases are primarily selected from airports, shopping malls, and stadiums because they, like stations, are also large infrastructures that deal with massive users.



\subsubsection{Literature}
\label{Literature}
Literature is searched through four themes. 17 papers are finally selected and used for providing heuristics for research-by-design (see \href{https://github.com/1309928130/Paper3_SupplementaryMaterials/blob/main/SupplementaryMaterials.pdf}{Supplementary Materials, fig. \ref{fig:sankeyLR}}). 

The first theme of searching is demand-responsive related concepts (Section \ref{Demand-responsive as an objective}), including ‘resilience,’ ‘flexible,’ `dynamic,' `adaptive,' ‘robust,’ `movable,' `reconfigurable,' `multi-functional,' ‘temporary,’ and so on (see \href{https://github.com/1309928130/Paper3_SupplementaryMaterials/blob/main/SupplementaryMaterials.pdf}{Supplementary Materials} for more explanations). The second theme is specifically for searching heuristics that deal with overcrowding: pedestrian simulations of railway stations. The third theme is specifically for searching heuristics that deal with emptiness: events (or activities) in public spaces (or urban/open spaces,’ ‘places,’ and ‘landscapes’). The fourth theme includes the station-city consideration: station-city integration from the urban design perspective. Google Scholar was used as the database for searching, considering its diverse range of literature types, making it suitable for design-related content beyond typical scientific knowledge.


\subsection{Data processing}
\label{Data processing}

\subsubsection{Examine relevance}
\label{General data processing}
We examine each case and literature paper regarding its relevance to design objectives and dimensions (See \href{https://github.com/1309928130/Paper3_SupplementaryMaterials/blob/main/SupplementaryMaterials.pdf}{Supplementary Materials}, fig. \ref{fig:Examine_relevance_dataprocess}, and Sections \ref{Design objectives}, \ref{Dimensions of spatial and managerial interventions}). 





\subsubsection{Comparing cases} 
We also compare various pairs of cases to see evident commons or differences.

\subsubsection{Research by design}
\label{Research by design}
Research by design (RbD) is a research method that uses design to develop knowledge. We adopt the definition of RbD proposed by \citet{dejongeWAYSStudyResearch2002}, in which RbD is defined as \textbf{changing the design itself or its context} and studying the \textbf{effects} of the transformations (Fig. \ref{fig:methodologyframework}).

More specifically, in this research, the changes/transformations we made include transferring strategies between different types of buildings/spaces; transferring strategies between different scales; testing with different shapes, sizes, positions, directions, and so on \citep{dadftArchitecturalFormDesign1991}; making novel combinations/syntheses of different objectives and measures. 

Fig. \ref{fig:researchbydesign} is an example of using RbD to study the relationship between event typology and the positioning of city passages. (See also  \href{https://github.com/1309928130/Paper3_SupplementaryMaterials/blob/main/SupplementaryMaterials.pdf}{Supplementary Materials} explaining the validity of research by design.)

\begin{figure}[h]
    \centering
    \includegraphics[width=\linewidth]{Paper3Imgs_pdfs/researchbydesign.pdf}
    \caption{An example of research by design}
    \label{fig:researchbydesign}
\end{figure} \FloatBarrier

RbD facilitates \textbf{new} knowledge production. New knowledge is urgently needed to extend the usefulness of existing knowledge in various situations in this research: when meaningful heuristics are found while no real-world solutions exist; when there is an existing solution and it has the potential to be extended; when some heuristics from literature is purely about activity and not related to spatial features, so spatial translation is needed. 

RbD facilitates our investigation of the \textbf{effects} of design interventions. It is common that the context and design solutions are not clear during the design process. `Research by design' helps designers to make both context and design interventions more explicit, hence making the effects of design interventions more evident to judge.

\subsection{Criteria for evaluation}
\label{Criteria for evaluation}
We evaluate the design principles with two criteria: 1) whether the findings meet design objectives (Section \ref{Design objectives}); 2) how practical they can be used in practice \citep{nijhuisDesignrelatedResearchLandscape2012}. Firstly, we check the design solutions by ourselves using a table where design objectives are listed against principles (See \href{https://github.com/1309928130/Paper3_SupplementaryMaterials/blob/main/SupplementaryMaterials.pdf}{Supplementary Materials}, fig. \ref{fig:evaluateDP}). Secondly, we do interviews with field experts to reflect on the design principles. Simulations \citep{chenEvaluateUserSatisfaction2024} are used for some design principles (such as principle 17) whose effects are significantly uncertain.

\section{Results}
\label{Results}
Our research resulted in two types of results: 1) the (over twenty) individual \textbf{design principles} and 2) the multiple \textbf{perspectives} to view these principles. In the following sections, we hierarchically present them for better clarity: first, we describe an overview of the principles, then introduce the multiple perspectives, and finally, the individual principles. These perspectives and principles can be checked on the interactive website we developed (\url{http://c1309928130.pythonanywhere.com/}).

\subsection{Overview of the design principles}
\label{Overview of the design principles}
Within the time frame of this research, more than twenty design principles are developed. It is evident that more potential design principles can be found in future research following the same methodology. Literature, cases, and research-by-design all provide substantial heuristics for the final result (fig. \ref{fig:sourcesofheuristics}). Different design principles have their varying strengths and limitations (\href{https://github.com/1309928130/Paper3_SupplementaryMaterials/blob/main/SupplementaryMaterials.pdf}{Supplementary Materials}, fig. \ref{fig:evaluateDP}).

\begin{figure}[h]
    \centering
    \includegraphics[width=\linewidth]{Paper3Imgs_pdfs/sourcesofheuristics.pdf}
    \caption{Sources of heuristics}
    \label{fig:sourcesofheuristics}
\end{figure} \FloatBarrier


\subsection{Perspectives to view the design principles}
\label{Perspectives to view the design principles}
During the research process, parallel with the new principles being developed, various perspectives emerged. These multiple perspectives can be seen as theoretical lenses through which to view the design principles (DPs) or typologies of the DPs. The multiple perspectives provide diverse ways for designers to relate to the individual principles we developed. The perspectives also guide designers in finding new principles that follow the same way of thinking.

\subsubsection{As patterns within a network}
\label{Patterns within a network}

We first arrange the resulting design principles in a \textbf{network} from the district level to the building level (fig. \ref{fig:Pattern_language_network}). This perspective is inspired by Alexander's pattern language \citep[p. 314]{alexanderTimelessWayBuilding1979}, whose structure is a network usually organized in a hierarchical \textbf{scalar} order. The cross-scale connections between principles are important because each principle relies on the smaller principles it contains, and sometimes also the larger principles in which it is contained. For example, to \hyperref[Connect with neighborhoods]{Connect with neighborhoods (8)} at the district level, station areas should properly \hyperref[(20 21) City passages and their positioning]{Position city passages (21)} and, design \hyperref[(20 21) City passages and their positioning]{City passages (20)} at the building level; Maybe \hyperref[(25) City axes]{City axes (25)} on a sub-city level are also needed. For another example, to generate leisure events, station areas can \hyperref[(16) Add installations and facilities]{Add installations and facilities (16)}, but the space where these installations and facilities are located should be \hyperref[(7) Human-oriented spaces]{Human-oriented spaces (7)}.

\begin{figure}[h]
    \centering
    \includegraphics[width=\linewidth]{Paper3Imgs_pdfs/Pattern_Network_8M.pdf}
    \caption{The design principles (patterns) within a network}
    \label{fig:Pattern_language_network}
\end{figure} \FloatBarrier

\subsubsection{As principles about different station components}
\label{Other perspectives}
The design principles can be categorized according to the component types of station areas which they deal with. On top of the aforementioned network, we already categorized the principles as related to \textbf{function, mass, space, flow \& path, mobility nodes}, and \textbf{leisure events} (fig. \ref{fig:Pattern_language_network}), most of which are common perspectives used in urban, architectural, and railway station design. 

As part of the built environment, station areas can be seen as function versus form, as the mass versus the void (space) \citep{rossiArchitectureCity1982, trancikFindingLostSpace1986},  as constellations of mobility nodes (including railway platforms, bus stops, tram stops, bike parking, car parking, and so on) \citep{vanderspekConnectorsWayTransferring2003}; Space users in station areas move along various paths generating pedestrian flow; Additionally, to be demand-responsive, station areas should also support leisure events during non-peak times when they are not used for transportation. 

\subsubsection{Other perspectives}
The design principles can be viewed from more perspectives (fig. \ref{fig:perspectivematrix}): From the problem-solving perspective -- addressing \textbf{overcrowding and emptiness} -- some principles address the former, some address the latter, and some address both. From a \textbf{spatial versus managerial} perspective, some principles are purely about spatial interventions, while others require managerial interventions; In this sense, principles can also be seen as \textbf{fixed and semi-flexible}.  



The principles can also be seen as related to \textbf{different event types} (see \href{https://github.com/1309928130/Paper3_SupplementaryMaterials/blob/main/SupplementaryMaterials.pdf}{Supplementary Materials}, fig. \ref{fig:eventtypology_principles}). For instance, if there are \textbf{large} events (e.g., mass gatherings and concerts), then design principles \hyperref[(11) Adaptive redundant spaces with path regulation]{Adaptive redundant spaces with path regulation (11)} and \hyperref[(7) Human-oriented spaces]{Human-oriented spaces (7)} should be adopted to reserve large open spaces that are event-supportive; While if the stations are designed in a scattered layout (to \hyperref[(P4 P5) Vibrate city environment or increase accommodating capacity by scattered mobility nodes]{Vibrate city environment by scattered mobility nodes (5)}), then possibly only \textbf{small} events can be held in station areas. 



The design principles can be viewed as rules of \textbf{elements and layouts} \citep{liuMappingLandscapeSpaces2020}; However, elements and layouts are actually similar unless a certain scale is chosen to analyze them. In the aforementioned network (fig. \ref{fig:Pattern_language_network}), the layout on a smaller level, can be the element on a higher level (e.g., \hyperref[(20 21) City passages and their positioning]{City passages (20)} can be parts of \hyperref[(25) City axes]{City axes (25)}). In fact, just like a layout, an element `is itself entirely a pattern of relationships' \citep[p. 89]{alexanderTimelessWayBuilding1979}; it defines relationships between different entities, possibly on different scalar levels (e.g., Tunnels can be called `city passages' only when they link different city areas, otherwise they are just meaningless floating cuboid geometries).

Still, more perspectives can be developed as long as they make sense to readers of this paper who want to use the design principles. For example: 1) naming a perspective `practice paradigms,' as in planning practice, there has been a paradigm shift from \textbf{car-oriented to human-oriented}  \citep{bertoliniStreetsTrafficStreets2020}, and two of our design principles (Principles \hyperref[(2) Suitable general layouts of the station and city]{Suitable general layouts of the station and city (2)}, \hyperref[(7) Human-oriented spaces]{Human-oriented spaces (7)}) are related to these paradigms. 2) Naming a perspective `\textbf{Bottleneck-related},' as several design principles (Principles \hyperref[(18 19) Redundant spaces or setting apart bottlenecks]{Set apart bottlenecks (19)}, \hyperref[(22 23) Reduce barriers to ease flow]{Reduce barriers to ease flow (22 \& 23)}) deal with the bottlenecks in the flow \& paths. 3) Naming a perspective `\textbf{Passage types},' as passages can be of different types according to the entities they connect, such as connecting city areas to city areas, city areas to stations, or inner station areas to inner station areas. 

\begin{figure}[h]
    \centering
    \includegraphics[width=0.9\linewidth]{Paper3Imgs_pdfs/perspectivematrix.pdf}
    \caption{The design principles viewed from different perspectives}
    \label{fig:perspectivematrix}
\end{figure} \FloatBarrier

In summary, with the above multiple perspectives, designers can view the design principles and relate them using various logic; and researchers/designers can develop new principles following the same ways of thinking.

\subsection{Design principles}
\label{Design principles}
In this section, we present each design principle in the form of \textbf{Context-Problems/Potentials-Solutions-Examples-Limitations(/Reflections)} whenever possible. All the principles are aimed at addressing the general overcrowding and/or emptiness problems, in the context of demand-responsive and station-city integration, as described in chapter \ref{Introduction} \hyperref[Introduction]{Introduction}. Some individual principles may aim at their own problems and contexts, which are detailed in the relevant description texts when necessary. Real-world examples are provided when available. For validation purposes, each principle is reviewed by field experts during expert interviews, and the comments are incorporated into the relevant description texts.


{    % Group to localize changes
    \renewcommand{\thesubsubsection}{}
    \makeatletter
    \renewcommand{\subsubsection}{\@startsection{subsubsection}{3}{\z@}%
      {-3.25ex\@plus -1ex \@minus -.2ex}%
      {1.5ex \@plus .2ex}%
      {\normalfont\normalsize\bfseries}}
    \makeatother



\subsubsection{(1) Flexible use}
\label{(1) Flexible use}

\textbf{Context \& Problems}: Stations can suffer from fluctuations of use, including overcrowding during peak times and emptiness during non-peak times (See \ref{Introduction} \hyperref[Introduction]{Introduction}).

% During peak times or when there are failures of railway lines, station areas can be overcrowded with detained passengers. 

\textbf{Solutions}: Flexible use is the solution that addresses fluctuation in use (fig. \ref{fig:flexible_use}). Two types of spaces can be flexibly used: 1) Stations' own spaces, and 2) City neighborhoods' spaces near the stations. The ideal scenarios are: During stations' peak times or emergent times, redundant spaces are used for accommodating passengers; during non-peak times, redundant spaces are used for city leisure activities or events (such as open markets, outdoor music festivals, and carnivals).

Stations need to \hyperref[Connect with neighborhoods]{Connect with neighborhoods (8)}, so that spaces of both cities' spaces can be flexibly used by stations accommodating peak-time passengers. For non-peak-time leisure use, it is necessary to further \hyperref[(22 23) Reduce barriers to ease flow]{Reduce barriers to ease flow (22) \& (23)}, provide \hyperref[(7) Human-oriented spaces]{Human-oriented spaces}, with leisure events and \hyperref[(6) Make events visible by aligning main paths and open spaces]{Make events visible by aligning main paths and open spaces (6)}.

Station areas need to also to be \hyperref[(12) Programming considering the time dimension]{Programmed considering the time dimension (12)} for flexible use of spaces.

Stations can be designed as \hyperref[(14 15) Flexible buildings or building components]{Flexible buildings or building components (14) \& (15)}, to have a responsive capacity for accommodating the peak time passengers. 

In some cases, railway stations are used temporarily, or temporarily drastically increase capacity, to facilitate city events or nearby stadiums, such as the Rotterdam Stadion Station (for Feyenoord Stadium), Victoria Station Manchester (for AO Arena), Zandvoort Station (for Formula 1 car racing events), Bijlmer Arena (for multiple nearby venues including Johan Cruyff Arena and Ziggo Dome). However, this flexible change of capacity is more of a pure managerial intervention rather than a spatial design intervention.

\begin{figure}[h]
    \centering
    \includegraphics[width=0.5\columnwidth]{Paper3Imgs_pdfs/index_pictures.027.pdf}
    \caption{Flexible use}
    \label{fig:flexible_use}
\end{figure} \FloatBarrier

\textbf{Concepts clarifications}: Of notice, in the `flexible use,' the `time dimension' is intrinsically considered, while some similar terms do not necessarily incorporate it. To better emphasize the \textbf{time dimension} nature, terms like `flexible,' `temporary,' `reconfigurable,' `adaptive,' and `changeable' are more appropriate than terms like `mixed-use,' `multi-functional,' `all-in-one,' and `sharing.' Because the latter can mean something has multiple functions working at the same time, not necessarily at different times. In the practice of urban planning and metro station development, similar to `flexible use,' the concept `mixed-use' \citep{rowleyMixeduseDevelopmentAmbiguous1996} has long been promoted. But `mixed-use' is not evident to readers as relates to time, even though a few scholars also incorporate the time dimension \citet{hoppenbrouwerMixeduseDevelopmentTheory2005} as the fourth dimension, beyond the 3D physical spatial dimensions, into the `mixed-used' definition. 

Flexible sharing with neighborhoods at least has two positive implications: 1. Given the same investment, station-city will have larger capacities. 2. Given the same capacities, smaller spaces are needed, which means less investment and less redundant space (and hence less emptiness during non-peak times).

\textbf{Limitations}: Flexible use has its limitations, as it can potentially add to the management efforts for stakeholders of stations and cities (e.g., see the chaos in Bijlmer Arena station when the management capacity is surpassed \citep{aseniyaCrowdControlBijlmer2019}), and as it can exacerbate overcrowding and emptiness if time dimension is not well considered.








\subsubsection{(2) Suitable general layouts of the station and city}
\label{(2) Suitable general layouts of the station and city} 

\textbf{Context \& Problems}: Stations and cities have different elements (and can be seen as the mass and the void); depending on the contexts, a suitable general layout of them can make the spatial allocation more understandable and rational. The following are some options of the general layouts: 
% - overcrowding
% - emptiness

\textbf{Solutions \& Examples}: The general layouts of stations' and cities' elements can be 
1) separated;
2) connected - e.g., Rotterdam Central Staiton and Utrecht Central Staiton, where stations are \hyperref[Connect with neighborhoods]{Connect with neighborhoods (8)} through \hyperref[(25) City axes]{City axes (25)}, and \hyperref[(20 21) City passages and their positioning]{(20) City passages and their positioning};
3) stacked - e.g., Montparnasse Station, upon which are apartment and offices buildings, and rooftop parks;
4) mixed - e.g., Shibuya Station and Shinjuku Station, where stations' spaces and cities' spaces are inseparable;
(fig. \ref{fig:generallayout}) \citep{shikataIntegratedStationCityDevelopment2013, qiAnalysisStationCitySynergetic2019, dewildeMultipleUseSpace2006}.

Whether mixing or separating stations' and cities' elements depends on the specific project contexts. Bringing non-transport city functions into the station areas, has several benefits, including intensifying land use, more commercial revenue, vibrating the spaces, and so on. However, this can potentially also cause conflicting use between the stations and cities, between passengers and non-passengers. Therefore, sometimes it is necessary to \hyperref[(3) Set apart non-transport function]{Set apart non-transport function (3)}.

These general layouts can be applied on different spatial scales, as long as the layouts facilitate designers' or stakeholders' understanding of the spatial configurations of station areas/buildings/elements.

\textbf{Limitations}: There are some limitations of trying to apply the general layouts: 1) The layouts can be meaningless or merely abstract concepts if they not related to the actual functioning of the stations' cities' components. 2) Also, sometimes defining stations and city elements is hard work itself. These elements may not have clear identifiable boundaries, or they have \hyperref[(1) Flexible use]{Flexible use (1)}.

\begin{figure}[h]
    \centering
    \includegraphics[width=0.5\columnwidth]{Paper3Imgs_pdfs/index_pictures.020.pdf}
    \caption{Different general layouts of the station and city}
    \label{fig:generallayout}
\end{figure} \FloatBarrier







\subsubsection{(3) Set apart non-transport function}
\label{(3) Set apart non-transport function}

\textbf{Contexts}: Setting apart non-transport function is one among several \hyperref[(2) Suitable general layouts of the station and city]{Suitable general layouts of the station and city (2)}. 

In the station-city integration practice, and the related and more popular topic of transit oriented development (TOD), mixed use is a fundamental principle. Despite the the mixed-use strategy have many benefits, in some situations, it also causes problems, and these problems can be easily overlooked.

\textbf{Problems}:
Mixing of stations with cities, if not done properly, can make the station buildings unrecognizable, and can bring interference in flow between passengers and non-passengers.

% in Section \ref{Spatial scales of concern: the district and building levels} \hyperref[Spatial scales of concern: the district and building levels]{Spatial scales of concern: the district and building levels}, 


\textbf{Solutions \& Example}:
Therefore, setting apart non-transport function is a solution in the above situation. One representative example is the Utrecht Central Station area. Before the latest redevelopment, the Utrecht Central Station building is highly mixed with the adjacent Hoog Catharijne shopping mall building. This mixture raised complaints from passengers as the station building is unrecognizable, passengers felt lost once they got off the train arriving this area, and the passages across the station building was always overcrowded with passengers and city users. During the latest development, a east well-designed plaza is added, separating the station building from the shopping mall building (fig. \ref{fig:utrechtsetapartnonstation}), and endowing the station a clear image identity. A spacious footbridge, with street-like qualities, is built alongside the station building for pedestrians to pass-by freely.

\textbf{Limitations}: The distinction between separation and mixing can be ambiguous, so apply this design principle (i.e., Setting apart non-transport function) still depends much on designers' experience.

\begin{figure}[h]
    \centering
    \includegraphics[width=0.5\columnwidth]{Paper3Imgs_pdfs/index_pictures.016.pdf}
    \caption{Set apart non-transport function}
    \label{fig:utrechtsetapartnonstation}
\end{figure} \FloatBarrier




\subsubsection{(4) \& (5) Vibrate city environment or increase accommodating capacity by scattered mobility nodes}
\label{(P4 P5) Vibrate city environment or increase accommodating capacity by scattered mobility nodes}

\textbf{Contexts \& Problems \& Solutions}: A layout of scattered mobility nodes can be applied for the following purpose:
1) To let railway passengers, who are potential customers of city stores, go through city districts, vibrating the city environment.
2) To increase the peak time capacity of the mobility sites -  passengers are more evenly dispersed to multiple sites - instead of concentrating on a limited site and, therefore, are less prone to overcrowding or even stampedes.
% increase capacity (Building-level, P5) (fig. ). Given the same volume of pedestrians, a more scattered layout of mobility sites enables less dense concentrations of pedestrians in each site. 

Scattered layouts of mobility nodes can be applied at the district level or the building level (figs. \ref{fig:districtscattered}, \ref{fig:buildingscattered}). 

\textbf{Examples}: Cases like Shanghai railway station adopt the layout of scattered mobility sites. 

\textbf{Limitations}: The side effects include: 1) that more mobility sites result in harder wayfinding, and 2) harder transferring between mobility sites (Therefore, if a station has a heavy transfer load between different mobilities, it is fine to use a scattered layout of mobility nodes; While if it has light transfer load between different mobilities, it is not appropriate to use this layout).

\begin{figure}[h]
    \centering
    \includegraphics[width=0.5\columnwidth]{Paper3Imgs_pdfs/index_pictures.017.pdf}
    \caption{Vibrate city environment by scattered mobility nodes}
    \label{fig:districtscattered}
\end{figure} \FloatBarrier
\begin{figure}[h]
    \centering
    \includegraphics[width=0.5\columnwidth]{Paper3Imgs_pdfs/index_pictures.014.pdf}
    \caption{Increase transport capacity by scattered mobility nodes}
    \label{fig:buildingscattered}
\end{figure} \FloatBarrier




\subsubsection{(6) Make events visible by aligning main paths and open spaces}
\label{(6) Make events visible by aligning main paths and open spaces}

\textbf{Contexts}: If a station suffers from emptiness, to vibrant the environment and reduce emptiness, leisure events can be organized.

\textbf{Problem}: But not every space is supportive of leisure events. If events happen in some places that are invisible to pedestrians (e.g., due to level changes, street-level sight lines are blocked) or away from the main path, it will be hard for events to attract enough people, and therefore, the empty station spaces remain empty. For instance, in fig. \ref{fig:pathwithevent}, if events are held in both stations, the environment in the right scenario will be energized significantly while the environment in the left will not.

\textbf{Solutions \& Examples}: The alignment between open spaces and event-goers' (leisure users') paths should be well-designed, to make sure events are visible, or even accessible for pedestrians to participate - depending on the event types. In the design stage, designers can either align spaces with paths or vice versa, aligning pedestrians' paths with spaces. 

Many different types of events can be held in station areas. For example, open markets, outdoor music festivals, carnivals, and so on \citep{smithStagingCityEvents2021}
\textbf{Activities} can be seen as small-scale events. 

The events' visibility can be enhanced by using \hyperref[(24) Stairs as stages or seats]{Stairs as stages or seats (24)}.

The event-supportive spaces are usually \hyperref[(7) Human-oriented spaces]{Human-oriented spaces (7)} that support \hyperref[(1) Flexible use]{Flexible use (1)}.

The paths that event-goers (leisure users) are willing to walk through should be considered in a large network of paths, which can consist of \hyperref[Connect with neighborhoods]{Connections with neighborhoods (8)} that \hyperref[(22 23) Reduce barriers to ease flow]{(22 23) Reduced barriers to ease flow}, and \hyperref[(20 21) City passages and their positioning]{City passages and their positioning (20) \& (21)}.

% \textbf{Limitations}: 

\begin{figure}[h]
    \centering
    \includegraphics[width=0.5\columnwidth]{Paper3Imgs_pdfs/index_pictures.001.pdf}
    \caption{Alignment between open spaces and main paths}
    \label{fig:pathwithevent}
\end{figure} \FloatBarrier







\subsubsection{(7) Human-oriented spaces}
\label{(7) Human-oriented spaces}

% - flexible use - non-transport activities; leisure events

\textbf{Contexts}: It is debatable if station areas should prioritize human uses, as stations are fundamental transportation infrastructures, and hence, people would think of stations as primarily for vehicles. This debate exists in the urban planning field for the general urban environment and for other transport infrastructures, including streets \citep{bertoliniStreetsTrafficStreets2020}.

In the past decades, there has been a paradigm shift from car-oriented to human-oriented planning; and for the design of railway station areas, at least in several situations, human-centered spaces should be a priority: 1) When station areas are shared with the cities and used by non-transport users (see principle \hyperref[(1) Flexible use]{Flexible use (1)}), or 2) when it is necessary to organize leisure events to reduce space emptiness (see chapter \ref{Introduction} \hyperref[Introduction]{Introduction}).

\textbf{Problems}: To support city non-transport use or leisure events, station spaces should be attractive. However, the human experience is often neglected in the design of some stations, where problematic features make them unattractive to people: 1) traffic elements dominating the spaces in railway station areas; 2) open spaces surrounded by traffic roads; 3) station buildings and spaces being monolithic, which makes them visually boring, and hard for people to move on foot for the long distance between facilities.

\textbf{Solutions \& Examples}: Human-oriented spaces should have good spatial qualities that enhance human experience, as many as possible. The qualities depend on different cultural contexts, project contexts, different users with different evolving demands, and so on. It is impossible to exhaustively define all the qualities, and that is perhaps why Alexander calls it the `quality without a name' \citep{alexanderTimelessWayBuilding1979}. Nevertheless, in urban design field, scholars have explicated some common qualities or spatial features, for instances: \textbf{human scale}; \textbf{enclosed} spaces with well-defined boundaries \citep{sitteCityPlanningAccording1889}; positives \textbf{edges} that promote social activities, urban life \citep{gehlLifeBuildings1987, whyteSocialLifeSmall2001}, and recognition/identity of different areas \citep{jacobsDeathLifeGreat1961, lynchImageCity1960}; \textbf{complexity/diversity} of visual elements \citep{ewingMeasuringUnmeasurableUrban2009}. 

The spatial layout, besides space qualities, also matters for human experience.  In fig. \ref{fig:humancenteredspace}-left, we try to illustrate some features regarding the layout: a human-centered space should be a coherent space at a core position in the station areas or the station building, free from traffic elements, and free from intensive pass-by transport flow so that people are possible to stay there. In contrast, the layout in fig. \ref{fig:humancenteredspace}-right does not have these features.

A significant example is the station plaza of the new Rotterdam Central Station, as against the old one (fig. \ref{fig:Rot_old_new}). In the old scenario, the core area was occupied by transport elements, including trams, buses, car parking, and a metro station entrance building; Walkable spaces are located around this core area; Hardly any area is possible for people to stay. In the new scenario, the core area is designed for humans; Metro entrances, tram platforms, and bus stations are moved aside; Car parking is located underground. During large events such as the Rotterdam Summer Carnival, the new station plaza is used without inference to the normal functioning of the station's transport operations. While in the old scenario, assuming there would have been an event, the old plaza was hardly possible to use unless all the transport elements stopped functioning.

Addressing high fluctuations of use towards demand-responsive station areas, human-oriented spaces can be supplemented with several smaller design principles, including \hyperref[(9 10) Smooth level changes by landscape design]{Smooth level changes by landscape design (9)}, \hyperref[(11) Adaptive redundant spaces with path regulation]{Redundant spaces with path regulation (11)}, \hyperref[(16) Add installations and facilities]{Add installations and facilities (16)} and \hyperref[(17 26) Reconfigurable elements or spaces]{Reconfigurable elements or spaces (17) \& (26)}.

Pedestrian paths, although usually linear elements for passing through, can be designed as human-oriented spaces if these paths are also supposed to let people stay. From this perspective, human-oriented spaces can be used as a strategy to support \hyperref[(25) City axes]{City axes (25)} and \hyperref[(22 23) Reduce barriers to ease flow]{Reduce barriers to ease flow (22) \& (23)}.

\textbf{Limitations}: There are some limitations of human-oriented spaces, including 1) the possible higher cost to create such spaces; 2) certain good qualities in one case may not necessarily be good in another case.

\begin{figure}[h]
    \centering
    \includegraphics[width=0.5\columnwidth]{Paper3Imgs_pdfs/index_pictures.002.pdf}
    \caption{Human-oriented spaces versus vehicle-oriented spaces}
    \label{fig:humancenteredspace}
\end{figure} \FloatBarrier
\begin{figure}[h]
    \centering
    \includegraphics[width=0.5\columnwidth]{Paper3Imgs_pdfs/index_pictures.030.pdf}
    \caption{The old and the new Rotterdam central stations}
    \label{fig:Rot_old_new}
\end{figure} \FloatBarrier



\subsubsection{(8) Connect with neighborhoods}
\label{Connect with neighborhoods}

\textbf{Contexts}: In the following context, stations should connect with their city neighborhoods:
% - lack of space, sharing to increase the capacity during peak times 
% - flexible use - promote movement, energize the environment
% - when city axes are already planned

% 1) If a station has high fluctuations of use, it should connect with its neighborhoods so that: During peak times or emergent, the station can use city spaces for high capacity accommodating strained passengers. During non-peak times, city users can use the redundant spaces in the stations. 

When stations have high fluctuation of use, and \hyperref[(1) Flexible use]{Flexible use (1)} with neighborhoods is needed, for the purpose of temporarily using the neighborhoods' spaces to accommodating stations' strained passengers during peak or emergent times; 

\textbf{Problems}: In contrast, if there are no such connections, problems will occur: the station spaces suffer from overcrowding during peak or emergent times.


\textbf{Solutions}: Connecting with the neighborhoods with stations (fig. \ref{fig:connection}), therefore, becomes a solution. Connections/linkages include pedestrian paths, footbridges/skywalks, tunnels, and so on. Passengers can use these connections when they have to. 

Connections are typically smaller parts of longer paths. They often cause sudden changes in the width of the paths, so it is necessary to ensure that they are not bottlenecks of flow. The connections should have enough capacity for passengers to pass during peak times.

The connections can basically serve for pass-through usage, whereas more preferably, the connections should be designed in a pleasant way that \hyperref[(22 23) Reduce barriers to ease flow]{Reduce barriers to ease flow (P22)} (Unless the stations do NOT want attractive non-transport city users).

% Causes of peak flow: the surge of passengers due to events, holidays, festivals …Bottlenecks should be reduced as much as possible.

Within the full pedestrian paths in which stations and city neighborhoods are connected, \hyperref[(20 21) City passages and their positioning]{City passages and their positioning (20) \& (21)} should also be considered.


% \textbf{Examples}:
% Rotterdam, Utrecht, negative examples: Beijing West, 
% \textbf{Limitations}:


\begin{figure}[h]
    \centering
    \includegraphics[width=0.5\linewidth]{Paper3Imgs_pdfs/index_pictures.010.pdf}
    \caption{Connect with neighborhoods}
    \label{fig:connection}
\end{figure} \FloatBarrier




\subsubsection{(9) \& (10) Smooth level changes by landscape design}
\label{(9 10) Smooth level changes by landscape design}

% - bottleneck adds to overcrowding
% - barrier of leisure flow - level changes that draw too much attention of pedestrians


\textbf{Context \& Problems}: Level changes are common in station areas due to many station buildings being multi-floor, or having topographic variations in the surrounding landscape. Using stairs naturally becomes a solution for level changes (see also design principle \hyperref[(24) Stairs as stages or seats]{Stairs as stages or seats (24)} can also be considered). However, walking on stairs is not a smooth experience - it requires careful attention from pedestrians in case of falling down, and it is even more dangerous during overcrowding. To further advance \hyperref[(7) Human-oriented spaces]{Human-oriented spaces (7)}, alternatives are needed.

\textbf{Solutions \& Examples}:
Using landscape design (figs. \ref{fig:jiaxing}, \ref{fig:arnhem}) helps to smooth the level changes, as compared to stairs. It would be useful, especially in flow-intensive areas, to prevent pedestrians from falling down. Landscape design can apply on both the district level (e.g., see fig. \ref{fig:jiaxing}, the plaza in the early design proposal of Jiaxing Railway Station, China, by MAD Architects) and the building level (e.g., the slop in the station hall of Rotterdam Central Station). When applied at the building level, it can even bring fresh spatial experience (e.g., Arnhem Central Station fig. \ref{fig:arnhem}).

\textbf{Limitations}: Landscape design smoothing level change can be expensive, especially at the building level. (The design of buildings like Arnhem Central Station is usually related to the so-called `parametric design,' which requires computer algorithms dealing with complex geometries with many engineering challenges). Simply using escalators could be a cheaper alternative.

landscape design smoothing level change may cause confusion in wayfinding for people who are not familiar with such level-changing environments.

\begin{figure}[h]
    \centering
    \includegraphics[width=0.5\columnwidth]{Paper3Imgs_pdfs/index_pictures.007.pdf}
    \caption{Smooth level changes by landscape design at the district level}
    \label{fig:jiaxing}
\end{figure} \FloatBarrier
\begin{figure}[h]
    \centering
    \includegraphics[width=0.5\columnwidth]{Paper3Imgs_pdfs/index_pictures.022.pdf}
    \caption{Smooth level changes by landscape design at the building level}
    \label{fig:arnhem}
\end{figure} \FloatBarrier





\subsubsection{(11) Adaptive redundant spaces with path regulation}
\label{(11) Adaptive redundant spaces with path regulation}

\textbf{Contexts \& Problems}:
Adaptive redundant spaces with path regulation should be used if a station has the following dilemma:
1) The station has \hyperref[(18 19) Redundant spaces or setting apart bottlenecks]{Redundant spaces or setting apart bottlenecks (18 19)} to reduce overcrowding during peak times; But
2) during non-peak times, the added long transfer distance by redundant spaces annoys passengers.

% Redundant spaces can make transfer distance longer if they are directly located between originations and destinations. So it seems a dilemma to retain redundant space for peak time use and ensure efficient transfer during non-peak time; 

\textbf{Solutions \& Examples}: Regulating paths solves the above peak-time redundancy and non-peak-time transfer efficiency dilemma. It is similar to \hyperref[(P13) Shortcuts or optimizing paths]{Shortcuts or optimizing paths (P13)}.

During peak times, it is a common practice to regulate the path by fences or lines (fig. \ref{fig:fencezigzag}) to ensure pedestrians move in an orderly queue, which prevents pedestrians from concentrating randomly, leading to stampedes. 
% Referring to this practice, during peak time, using fence-like measures can also regulate the path to make people walk through the redundant spaces, where they do not go during non-peak time, to avoid overcrowding. 

Instead of fences or lines, which are quite intrusive elements, paths can also be regulated by more user-friendly, even interactive design elements, such as urban furniture, \hyperref[(17 26) Reconfigurable elements or spaces]{Reconfigurable elements (17)} - water bodies, fountains, potted plants, art installations, stalls, and so on (fig. \ref{fig:Path_regula_reconfig_elements}).

To further satisfy the waiting passengers in the regulated zigzag paths, it is beneficial to create a pleasant environment of staying, compensating for the longer waiting time than usual.

During non-peak times, the path-regulating elements can be removed to let pedestrians move freely. \hyperref[(7) Human-oriented spaces]{Human-oriented spaces (7)} is a prerequisite for redundant space to be adaptively used during non-peak times.

In Chinese station cases, path regulations (using fences) are more commonly seen during peak times. In the Netherlands, however, it is less common. 

\textbf{Limitations}: Path regulations can easily raise complaints from passengers or nearby citizens; therefore, they should be used sparingly with strong necessity. 

\begin{figure}[h]
    \centering
    \includegraphics[width=0.5\columnwidth]{Paper3Imgs_pdfs/index_pictures.012.pdf}
    \caption{Adaptive redundant spaces with path regulation}
    \label{fig:fencezigzag}
\end{figure} \FloatBarrier
\begin{figure}[h]
    \centering
    \includegraphics[width=0.5\columnwidth]{Paper3Imgs_pdfs/index_pictures.019.pdf}
    \caption{Path regulation using reconfigurable elements}
    \label{fig:Path_regula_reconfig_elements}
\end{figure} \FloatBarrier





\subsubsection{(12) Programming considering the time dimension}
\label{(12) Programming considering the time dimension}

\textbf{Contexts \& Problems}: To reduce overcrowding and emptiness in stations, \hyperref[(1) Flexible use]{Flexible use (1)} is needed, and one planning strategy is programming station areas considering the time dimension (fig. \ref{fig:programmingsharing}).

\textbf{Solutions \& Examples}: Green parks and pedestrian-friendly streets, if programmed next to stations, can be temporarily restrictedly used by stations exclusively - for accommodating waiting passengers during railway peak times or for evacuation during emergent times. 

Similarly, large infrastructures such as sports stadiums, and music venues, since they are usually not used on a daily basis, can be shared with stations for railway peak time use. For example, in Guangzhou, China, multiple sports stadiums were planned for backup use of railway stations if any emergency occurred during Chinese New Year times \citep{xinhuanewsagencyGuangzhouSpringFestival2016}. This makes sense in this culture, as during the Chinese New Year, there are surges of railway passengers \citep{mccarthyChineseNewYear2018} who want to go to their hometowns for family gatherings, while at the same time, sports stadiums are mostly empty.

During non-peak times, the city spaces and infrastructures can be used as normal for their original non-transport leisure activities. Or, they can even take advantage of the railway as a more sustainable means of transportation \citep{newmanSustainabilityCitiesOvercoming1999}. For example, in Amsterdam, the Netherlands, several stadiums are planned next to Bijlmer Arena railway station, promoting the use of the railway. 


\begin{figure}[h]
    \centering
    \includegraphics[width=0.5\columnwidth]{Paper3Imgs_pdfs/index_pictures.018.pdf}
    \caption{Programming considering the time dimension (for flexible use)}
    \label{fig:programmingsharing}
\end{figure} \FloatBarrier

\textbf{Limitations}: Programming has its limitations for events with frequency from multi-weekly to yearly (which are our focuses in this research). We discovered only limited programming strategies for this kind of event. In contrast, for more frequent events such as commuting (with a frequency between daily to weekly, and is an issue more for metro stations instead of railway stations), there are common programming strategies that have long been established, such as mixing housing/hospitality with offices.






\subsubsection{(13) Shortcuts or optimizing paths}
\label{(P13) Shortcuts or optimizing paths}

% (This design principle is more suitable for large Chinese stations)
\textbf{Context \& Problems}: In some stations, the transfer between different mobilities is effect-consuming. For instance, at Beijing West Railway Station, if a passenger comes out of the metro and wants to go to the railway, he needs to check the ticket out of the metro, go the lengthy paths with level changes, and then check in at the railway with a strict security check, and a ticket check (fig. \ref{fig:buildshortcut}). Therefore, shortcuts in the paths or optimizing paths are needed. 

\textbf{Solutions \& Examples}: If there are shortcuts (direct links between the metro and the railway) in the paths (fig. \ref{fig:buildshortcut}), the passenger can save transfer distance and be relieved from double security checks (He already had a security check when he entered the metro system). For example, Beijing South Railway Station (RS), as a counterpart of Beijing West RS, provides direct transfer links between the metro and the railway. 

% If there is a shortcut enabling him to directly go from the metro station to the railway station (e.g., refer to Beijing South Station), his transfer experience (in terms of transfer time and distance) will be improved, and there will be less redundant space (hence overcrowding during peak times and emptiness during non-peak times) needed for passengers like him.

The shortcuts can be applied during non-peak times for transfer speed, while shut down during peak times for safety (as a lengthy path has more redundant spaces to accommodate the strained passengers, preventing overcrowding).

Shortcuts in paths are similar to \hyperref[(11) Adaptive redundant spaces with path regulation]{Path regulation for adaptive redundant spaces}. Their difference (in this research) is that shortcuts are more for (indoor) buildings, and path regulation is more for outdoor spaces (fig. \ref{fig:Pattern_language_network}).

\textbf{Limitations}: However, adding shortcuts means adding extra ingress/egress paths, in which extra electronic systems are needed, hence increasing costs for railway companies. It also adds more complexity for the system and for the users. The well-designed \hyperref[(7) Human-oriented spaces]{Human-oriented spaces (7)} along the normal paths will be missed if passengers choose the shortcuts over the normal paths.

\begin{figure}[h]
    \centering
    \includegraphics[width=0.5\columnwidth]{Paper3Imgs_pdfs/index_pictures.006.pdf}
    \caption{Shortcuts or optimizing paths}
    \label{fig:buildshortcut}
\end{figure} \FloatBarrier






\subsubsection{(14) \& (15) Flexible buildings or building components}
\label{(14 15) Flexible buildings or building components}

\textbf{Contexts \& Problems \& Solutions}:
If a station has huge fluctuations of passengers, to support \hyperref[(1) Flexible use]{Flexible use (1)}, they can be designed as a flexible building, or a building with some flexible components (including, such as the flexible gate machines) (figs. \ref{fig:theshed}, \ref{fig:Changeable_building_components}). 

Similar to the word `flexible,' temporarily used buildings solve the emptiness issue from the root. There are various types of buildings titled with concepts that enable a certain degree of flexibility. These concepts include `changeable,' `movable,' or `temporary buildings,' `flexible,' `multi-functional' (in temporal dimension), `reconfigurable,' `kinetic' \citep{parametricarchitectureWhatKineticArchitecture2024}, and so on \citep{geraedtsSuccessFailureFlexible2011}. These buildings have flexibility either by changing their use (programming) during different times, or their existence (assemblage) during certain periods. 

\textbf{Examples}:
Flexible buildings in the real world are not common due to the high cost and high maintenance. Here are a few non-station examples: 1) The Shed, which has a movable hall as an event venue. 2) Stadium 974 , which is a stadium made of containers and can be assembled and shipped (see also design proposal by \citet{smithShipStadiumTransforming2023}). For station cases, there are examples such as 3) temporary station buildings during the construction of the main new station buildings (e.g., in the Rotterdam Central Station case). In Shanghai South Railway Station and Beijing West Stations, many redundant checking machines are allocated in the gate areas, which can be used during peak times for higher capacity.

\textbf{Limitations}: Flexible buildings are expensive to build and maintain; They are like a fantasy against reality in most situations.

\begin{figure}[h]
    \centering
    \includegraphics[width=0.5\columnwidth]{Paper3Imgs_pdfs/index_pictures.008.pdf}
    \caption{A changeable building - the Shed}
    \label{fig:theshed}
\end{figure} \FloatBarrier
\begin{figure}[h]
    \centering
    \includegraphics[width=0.5\columnwidth]{Paper3Imgs_pdfs/index_pictures.024.pdf}
    \caption{Changeable building components}
    \label{fig:Changeable_building_components}
\end{figure} \FloatBarrier




\subsubsection{(16) Add installations and facilities}
\label{(16) Add installations and facilities}

\textbf{Contexts \& Problems}: In the following contexts, it will be helpful to add installations and facilities to vibrant the environment: 
1) If a station suffers from emptiness during non-peak times; and 
2) the station spaces are huge and boring.
3) The station spaces are human-oriented spaces that are supportive of events leisure activities or big leisure events

\textbf{Solutions}:
Adding installations and facilities (fig. \ref{fig:antwerpferrywheel}) is helpful to reduce emptiness by 1) stimulating events; 2) reducing the size of huge empty spaces to make them more human-scale; 3) making the scenes visually more diversified to have enough complexity. Installations can be interactive (e.g., ferry wheels) or just visually interesting (big inflatable art installations \citep{iversStagingUrbanLandscapes2018}).

The installations and facilities can also be \hyperref[(17 26) Reconfigurable elements or spaces]{Reconfigurable elements (17)} that can be rearranged. 

\textbf{Examples}:
For example, At Antwerpen Central Station, a ferry wheel is sometimes installed in the front square. At Rotterdam Central Station, different art installations are installed from time to time, on the front square and in the station hall.

\textbf{Limitations}: If the installations and facilities are not attractive by themselves, then they can deteriorate the space quality.
For instance, in the Antwerpen Central Station case, when the ferry wheel is installed, the corresponding facilities, such as ticket-selling points, can make the square messy. On the street level, pedestrians' views are much blocked by these facilities, and the beautiful facade of the station is less visible to pedestrians. Similarly, on the South Square of Beijing West Railway Station, many staircases, soldier standing points, the roofs of sunken atriums, police car parking points, and so on, also block pedestrians' views.

\begin{figure}[h]
    \centering
    \includegraphics[width=0.5\columnwidth]{Paper3Imgs_pdfs/index_pictures.003.pdf}
    \caption{Add installations and facilities}
    \label{fig:antwerpferrywheel}
\end{figure} \FloatBarrier




\subsubsection{(17) \& (26) Reconfigurable elements or spaces}
\label{(17 26) Reconfigurable elements or spaces}

\textbf{Contexts \& Problems}: Reconfigurable elements or spaces can be used for the following purposes:
1) To vibrate the environment if station spaces are empty during non-peak times.
2) To \hyperref[(11) Adaptive redundant spaces with path regulation]{Regulate path (17)} during peak times.
3) To further develop \hyperref[(7) Human-oriented spaces]{Human-oriented spaces (7)}, making these spaces more diverse with elements.
4) To make \hyperref[(16) Add installations and facilities]{Installations and facilities (16)} more flexible

\textbf{Solutions \& Examples}:
Reconfigurable elements are movable. They can be rearranged in different layouts for different use scenarios. They can be urban furniture (fig. \ref{fig:urbanfurniture}), potted plants, art installations, stalls, and so on. Reconfiguration of elements can be a spontaneous activity of users, such as people who use urban furniture will move the furniture by themselves.

Reconfigurable spaces can be set as different venues for different activities or events
(fig. \ref{fig:Reconfigurable_spaces}). For example, Bryant Park in New York is constantly reconfigured as a music venue, tennis courts, a grassland, and an ice skating rink \citep{iversStagingUrbanLandscapes2018}. 

Reconfigurable spaces or elements are similar to \hyperref[(14 15) Flexible buildings or building components]{Flexible buildings or building components (14) \& (15)}. Their difference is that the former is about outdoor spaces, and the latter is about (indoor) buildings.

\begin{figure}[h]
    \centering
    \includegraphics[width=0.5\columnwidth]{Paper3Imgs_pdfs/index_pictures.004.pdf}
    \caption{Reconfigurable elements}
    \label{fig:urbanfurniture}
\end{figure} \FloatBarrier
\begin{figure}[h]
    \centering
    \includegraphics[width=0.5\columnwidth]{Paper3Imgs_pdfs/index_pictures.026.pdf}
    \caption{Reconfigurable spaces}
    \label{fig:Reconfigurable_spaces}
\end{figure} \FloatBarrier


\subsubsection{(18) \& (19) Redundant spaces or setting apart bottlenecks}
\label{(18 19) Redundant spaces or setting apart bottlenecks}

\textbf{Contexts \& Problems}: In the following situation:
Pedestrians move between different mobility sites in station areas, from originations to destinations, and there may be bottlenecks(Section \ref{Viewing stations from the pedestrian flow perspective}). 
1) If a station has surges of passengers during peak times; and 
2) if the mobility sites or bottlenecks are too close, then the strained passengers will quickly be added up to overcrowding.
Therefore, it is necessary to add redundant spaces or, equivalently, set apart bottlenecks.

\textbf{Solutions}: Redundant spaces (between originations to destinations, between different mobility sites, or between bottlenecks) reduce overcrowding in two ways:
1) Redundant spaces make pedestrians walk longer. Since different pedestrians walk at different speeds, the surges of pedestrians at the originations will be smoother when the pedestrians arrive at the destinations.
2) Redundant spaces are extra spaces that accommodate strained pedestrians.
Redundant spaces can be set at the district level (fig. \ref{fig:layout_setapart}) as well as at the building level (fig. \ref{fig:build_setapart}).

\textbf{Examples}: One example of redundant spaces being used is the Allianz Arena Munich. From the nearest train station to the stadium, it takes over 1km distance to walk. The massive space in between is redundant space that can accommodate a large number of people.
% Setting apart bottlenecks is similar to adding redundant space in between. If two bottlenecks are located too close, detained pedestrians caused by each of them are further concentrated, making the situation worse. Setting the two apart allows more (redundant) space and capacity to accommodate the aggregation of detained pedestrians.

\textbf{Limitations}: Redundant spaces have the following limitations:
1) Despite redundant spaces being necessary for safety concerns (i.e., prevent overcrowding) during peak times, the long distances added by the redundant spaces can be unacceptable for passengers during non-peak times. In this situation, temporary \hyperref[(P13) Shortcuts or optimizing paths]{Shortcuts or optimizing paths (P13)}, should be adopted, or make the spaces as \hyperref[(11) Adaptive redundant spaces with path regulation]{Adaptive redundant spaces with path regulation (11)}.
2) Overcrowding (or the more extreme - stampedes) is a tricky issue that is associated with many causes that are not yet fully understood by existing scientific research \citep{almeidaHumanStampedesUpdated2019}. Redundant spaces may not fully function if the crowd of people, due to crowd psychology, still concentrates in front of the destination areas (despite there being redundant spaces outside destination areas). In this situation, on-site management regulations are still needed to disperse people more evenly across the spaces.

\begin{figure}[h]
    \centering
    \includegraphics[width=0.5\columnwidth]{Paper3Imgs_pdfs/index_pictures.011.pdf}
    \caption{Redundant spaces or setting apart bottlenecks at the district level}
    \label{fig:layout_setapart}
\end{figure} \FloatBarrier
\begin{figure}[h]
    \centering
    \includegraphics[width=0.5\columnwidth]{Paper3Imgs_pdfs/index_pictures.015.pdf}
    \caption{Redundant spaces or setting apart bottlenecks at the building level}
    \label{fig:build_setapart}
\end{figure} \FloatBarrier







\subsubsection{(20) \& (21) City passages and their positioning}
\label{(20 21) City passages and their positioning}

\textbf{Context \& Problems}: It is common for railway tracks to cut city areas apart. City passages and their positioning should be designed/planned for the following purposes:
1) Facilitating passengers to move from one side of the station to the other side of the station.
2) Facilitating city users passing through the station areas that were previously blocked by railway tracks. The station areas will be more vibrant with more users passing by.
3) When stations are \hyperref[Connect with neighborhoods]{Connected with neighborhoods (8)}, the city passages across railway tracks, help to form coherent pedestrian paths at the district level, or \hyperref[(25) City axes]{City axes (25)} on an even larger scale. These paths or city axes facilitate pedestrian movement.

\textbf{Solutions \& examples}: At the district level, city passages across railway tracks should be positioned according to the city's neighborhood conditions (fig. \ref{fig:layoutpassage}). City passages may link the most dense neighborhoods/communities/city areas to maximize their usage. City passages should avoid bringing interference between transport passengers and city leisure users. If the passages are mainly for leisure users to pass by, then their positions should avoid areas where passengers intensively move. If the passages are mainly for transport passengers, their position should avoid areas where leisure users would stay or gather.

% enabling pedestrians to move freely across tracks or station buildings. This may, therefore, increase the usage of the spaces and facilities in station areas, which reduces emptiness during non-peak time. 
% (and event typology?)

If there are areas reserved for leisure events, to make events more visible or even interactive for pedestrians, the paths should be aligned close to these areas (see principle \hyperref[(6) Make events visible by aligning main paths and open spaces]{Make events visible by aligning main paths and open spaces (6)}).

At the building level, city passages can be designed with fine details, with enough lighting, and not too dark to make people feel unsafe (fig. \ref{fig:buildpassage}). Beyond city passages that link different city areas across railway tracks (city-to-city), there can be other types of passages that link station to city and station inner areas to station inner areas (e.g., the internal footbridge on platforms at the Railway Central Station) These passages also provide convenient movement and ease transfer.

Passages are commonly seen in Dutch railway stations, such as Amsterdam Central Station, Utrecht Central Station, and Rotterdam Central Station.

\textbf{Limitations}: The more passages, the more choices for pedestrians to pass through, but it takes financial cost to build passages. 
at the building level, it is often quite difficult to vertically arrange the passages with the railway platforms, the ground, and the internal passages that passengers move from platform to platform.

\begin{figure}[h]
    \centering
    \includegraphics[width=0.5\columnwidth]{Paper3Imgs_pdfs/index_pictures.023.pdf}
    \caption{Position city passages (at the district level)}
    \label{fig:layoutpassage}
\end{figure} \FloatBarrier
\begin{figure}[h]
    \centering
    \includegraphics[width=0.5\columnwidth]{Paper3Imgs_pdfs/index_pictures.013.pdf}
    \caption{The city passages (at the building level)}
    \label{fig:buildpassage}
\end{figure} \FloatBarrier






\subsubsection{(22) \& (23) Reduce barriers to ease flow}
\label{(22 23) Reduce barriers to ease flow}

\textbf{Contexts}: Barriers to flow should be reduced, especially in the following contexts: 
1) When stations suffer from emptiness, there is a need to promote the flow of people (leisure users from the city neighborhoods), letting them use the redundant station spaces for public life, leisure activities, events, and so on.
2) When \hyperref[(25) City axes]{City axes (25)} (in which the stations are located) are planned to promote the flow.
3) When stations already are \hyperref[Connect with neighborhoods]{Connected with city neighborhoods (8)}, but due to the connections' qualities being poor, only passengers driven by strong motives would pass through the connections, while leisure users are not willing to use the connections.

\textbf{Problems}: Barriers of flow are usually in the following forms with respective reasons: 
1) Stairs as barriers. Pedestrians need to make efforts to overcome level changes, and much attention is needed to walk on the steps in case of falling down.
2) Lengthy, narrow, and boring paths.
3) Tunnels, skywalks/footbridges - Pedestrians need to make efforts to overcome level changes - Pedestrians would rather walk on the ground if possible. Tunnels are usually quite dark, making pedestrians feel unsafe. Skywalks/footbridges are unpleasant to walk on if they are non-weather-protected, without shading from trees or building canopies.
4) Walking areas that are interfered with by intensive vehicle traffic - Pedestrians feel unsafe to walk through, and much attention is needed to prevent being hit by a car.
5) Strict security checking points (Common in Chinese stations) - Too much effort is needed to pass through them. 
6) Ticket gates - They make gated areas inaccessible for people without tickets.

Barriers can lead to commercial stores in station areas being shut down during non-peak times, since city leisure users, as potential customers, are not willing to overcome barriers to explore the stores.


\textbf{Solutions}: Therefore, it is necessary to Reduce barriers to ease flow, at the district level (fig. \ref{fig:beijingwestnorth}) and at the building level (fig. \ref{fig:buildflowcommercial}). The following interventions can be taken, depending on the barrier types:
1) Use ground-level paths as much as possible, to avoid level changes, and to continue the street-level qualities of the overall environment.
2) \hyperref[(9 10) Smooth level changes by landscape design]{Smooth level changes by landscape design (9) \& (10)}. (This is particularly useful with small level changes; For large level changes, using escalators is usually more appropriate).
3) For lengthy boring paths, design them like \hyperref[(7) Human-oriented spaces]{Human-oriented spaces (7)}, make them more pleasant with leisure facilities, otherwise use \hyperref[(P13) Shortcuts or optimizing paths]{Shortcuts or optimizing paths (P13)}
4) If footbridges/skywalks are adopted, create street-like qualities, such as weather-protected shading and convenience facilities.
5) Separate traffic roads from pedestrian paths, e.g., horizontally make alignment between the pedestrian flow and traffic flow to avoid overlapping, or vertically submerging traffic roads underground or elevating traffic roads.
6) Security checks and ticket gates are primarily management issues. For spatial design considerations, possibly \hyperref[(P13) Shortcuts or optimizing paths]{Shortcuts or optimizing paths (P13)} can help. 
7) Consider \hyperref[(20 21) City passages and their positioning]{City passages and their positioning (20) \& (21)} if station areas are cut apart by railway tracks.

% During non-peak times, city neighborhoods use the stations' spaces for non-transport public activities. Commercial services within the station areas will be more resilient as not only the passengers but also the city users are their potential customers.
% brings another benefit - increasing the use of commercial facilities. Given the same types and amount of commercial supply, spaces with no bottlenecks or barriers better support passengers to explore the whole variety of commercial facilities, and the commercial stores are more resilient the the local changes of users since they are patronized by passengers from the global network. 

% This reduces the redundant spaces needed, hence reducing and overcrowding as well as emptiness. 

% \textbf{Examples}:

\textbf{Limitations}: Given barriers are reduced, the flow of leisure users should not cause severe interference to passenger flow. Also, every measure has a (financial) cost, hence reducing barriers should bring enough benefits.

\begin{figure}[h]
    \centering
    \includegraphics[width=0.5\columnwidth]{Paper3Imgs_pdfs/index_pictures.009.pdf}
    \caption{Reduce barriers to ease flow at the district level}
    \label{fig:beijingwestnorth}
\end{figure} \FloatBarrier
\begin{figure}[h]
    \centering
    \includegraphics[width=0.5\columnwidth]{Paper3Imgs_pdfs/index_pictures.021.pdf}
    \caption{Reduce barriers to ease flow at the building level}
    \label{fig:buildflowcommercial}
\end{figure} \FloatBarrier







\subsubsection{(24) Stairs as stages or seats}
\label{(24) Stairs as stages or seats}

\textbf{Contexts \& Problems}: Stairs are common building components that primarily address level changes. Some stairs also serve for symbolic functions, for example, using stairs to elevate buildings (or building entrances) to make the buildings look more grand.

Typically, in station areas, stairs should be avoided as much as possible because they hinder the flow of pedestrians and may even result in people falling down. This issue is especially problematic if the flow is intensive. To address small level changes, instead of using stairs, it is better to \hyperref[(9 10) Smooth level changes by landscape design]{Smooth level changes by landscape design (9) \& (10)}. To address large level changes, instead of using stairs, it is better to use escalators. Given the same width, escalators have a larger capacity than stairs, and escalators, if not running, can still be used as stairs. 

However, in some situations, large stairs still exist in station areas: Those designed as backups of escalators for evacuations - like the grand stairs at the west of Utrecht Central Station. Or, stairs just happen to exist in the station areas - like the grand stairs of the Cologne Cathedral in front of the Cologne Central Station, and the indoor stairs in the hall of Antwerpen Central Station.

Stairs should be designed as potential stages for event performances or seats for the audience if the following two conditions are met:
1) Grand stairs, for whatever reasons, exist in railway station areas
2) Leisure events are supposed to be held in the station areas, as a measure to reduce emptiness.

% \hyperref[(6) Make events visible by aligning main paths and open spaces]{Make events visible by aligning main paths and open spaces (6)})

\textbf{Solutions}:
Grand stairs (in non-flow-intensive areas) can be used as audience seat areas, or as stages for performances during events. This takes advantage of the level change to leverage the see-and-be-seen relationship between audiences and performers in events. Such usage has been seen in various real-world cases: The grand stairs at the west side of Utrecht Central Station used during open-air exhibitions (fig. \ref{fig:utrechtstair}); The grand stairs in the hall of Antwerpen Central Station used for music festivals; And the grand stairs in front of Cologne Central Station used as stages for demonstrations.

\textbf{Limitations}: Stairs are typically unfavorable elements in the station areas, as mentioned in the above Contexts \& Problems text. If stairs are indeed used for events, the frequency of use is still questionable.




\begin{figure}[h]
    \centering
    \includegraphics[width=0.5\columnwidth]{Paper3Imgs_pdfs/index_pictures.005.pdf}
    \caption{Stairs as stages or seats}
    \label{fig:utrechtstair}
\end{figure} \FloatBarrier




\subsubsection{(25) City axes}
\label{(25) City axes}


\textbf{Context \& Problems}: If station areas suffer from emptiness, city axes, which are a common planning strategy in urban planning and urban design fields, can be planned. 


\textbf{Solutions}: City axes help to emphasize the connections, concentrate the flow of leisure users, and energize the environment.

City axes make the spatial structure of certain areas recognizable to decision-makers, thereby also promoting collaborations for different stakeholders to align the important spaces and facilities, for creating memorable experiences for users when they move in the environment.

City axes should be supported by smaller design principles to ensure pedestrians can really move within them and experience the environment. These smaller principles include: (stations) \hyperref[Connect with neighborhoods]{Connect with neighborhoods (8)}, \hyperref[(22 23) Reduce barriers to ease flow]{Reduce barriers to ease flow (22) \& (23)}, and properly design/plan \hyperref[(20 21) City passages and their positioning]{City passages and their positioning (20) \& (21)}. 

\textbf{Examples}: City axes have been successfully used in (re)development projects of some station areas, for example, the Utrecht Central Station area (fig. \ref{fig:City axes}), the Rotterdam Central Station area, and the Delft Station area.

\textbf{Limitation}: Of notice, city axes usually should be space-based (e.g., the Paris Historical Axis) instead of building-based (e.g., the axis of the Forbidden City in Beijing). Building-based city axes only work symbolically to demonstrate order, power, or spiritual values, and can hardly contribute to real space usage because people cannot move on the axes. 

\begin{figure}[h]
    \centering
    \includegraphics[width=0.5\columnwidth]{Paper3Imgs_pdfs/index_pictures.025.pdf}
    \caption{Two city axes were planned during the redevelopment of the Utrecht Central Station area}
    \label{fig:City axes}
\end{figure} \FloatBarrier

% \subsubsection{(26) Reconfigurable spaces}
% \label{(26) Reconfigurable spaces}

% adaptive redundant spaces VS path regulation
% networks of elements
% networks of point-of-interests
% networks of paths


% \subsubsection{(27) Event-supportive spaces}
% \label{(27) Event-supportive spaces}

% \subsubsection{(26) Nature solutions}
% \label{(26) Nature solutions}

% \subsubsection{(27) Compensate indoor space with more (outdoor) gray space}
% \label{(27) Compensate indoor space with more (outdoor) gray space}



}






\section{Discussion}
\label{Discussion}



\subsection{Contributions}
\label{Contribution}

This paper has two contributions to the design knowledge body: First, \textbf{contently}, it provides more than two dozen design principles that facilitate designers to address overcrowding and emptiness issues of station areas. This paper also provides multiple associated perspectives, which can be used to view the principles in different ways, and be used as theoretical lenses to help designers investigate new principles. The research outcome also enriches the content of relevant academic discussions, such as urban open space, responsive landscape, with a ‘loose-fit' feature \citep{thompsonUrbanOpenSpace2002}; and traffic space is public space \citep{bendiksTrafficSpacePublic2020}.

Second, \textbf{methodologically}, this paper provides a seven-step methodological framework, which can be generalized to develop other types of design principles. This methodological framework has its only values as compared to similar frameworks/workflows in literature. For example, compared to Christopher Alexander's way of developing his pattern language - through over eight years of project practice \citep{alexanderPatternLanguageTowns1977}, evaluate principles by feelings - our framework is more practical for much shorter time frames thanks to multiple data sources being used, and produce more objective outcomes thanks to expert reflections. Compared to the way \citet{vanderhoevenSignificanceUndergroundExperience2016} developed their design principles, our framework is a new alternative.

% The multiple established perspectives can serve as renewed theoretical lenses to explore more design principles in the future. There is no end to the number of design principles. Especially when it comes to public space user experience, as long as it can promote user satisfaction and perception, it can basically promote space use, thereby reducing emptiness. Weakening negative space use experiences such as wayfinding can reduce overcrowding, which will also include many wayfinding design principles. So this article cannot be too ambitious. The workflow ...





\subsection{Non-design interventions}
\label{Non-design interventions}
Many purely managerial (non-design) interventions are discovered during the research process. Some of these interventions can be more effective than spatial interventions, some can be more costly, and some can be fantasies at the current time. Providing as a takeaway knowledge, We report the following managerial interventions, which can be used to avoid the peaks of passenger flow, or to avoid overcrowding/stampedes:

Managerial interventions: 1) Reschedule holidays for the public; 2) reduce physical security checks and apply more AI and data technology for pre-assessing passengers. 3) Use city terminals - like those for the airport - to modify the logistics of long-distance travelers' packages and relieve the security checks at the station buildings; 4) manage regional-level transportation (e.g., there is a coordination of scheduling trains between several railway stations in Amsterdam city when there are sports events near Bijlmer Arena station, preventing people from coming to the stations when the station capacity is surpassed); 5) shutdown stations; 6) when organizing city-level events, use multiple scattered event spots - so that there will not be an over-concentration of massive people at one spot during special times such as when it is raining; 7) manage user expectations - notifying passengers that during peak time there will be longer transfer time than usual.

% - Regional level management of transportation … 
% - Instead of regulating paths directly near station buildings (such as on the station squares), it is also possible to regulate paths at a larger district level, (See the Solutions \& Examples in principle \hyperref[(12) Programming considering the time dimension]{Programming considering the time dimension (12)}), or even at a 
% - Sharing with other infrastructures. 



\subsection{Station typology implied by the cases}
\label{Station typology implied by the cases}

Considering that station is the type of building with the most variety \citep[p.~52]{richardsRailwayStationSocial1986}), a typology is usually of concern in studies of stations (e.g., see \citep{petersTrainStationArea2012}), and meanwhile, a universal typology is hardly achievable. In this research, we take two factors that are most fluctuation-related for case selection - fluctuation contexts and performance addressing fluctuation (see Section \hyperref[Cases]{Cases}). With the selected cases' other features listed bellow, the cases imply a certain station typology. Through these features, we can better tell how generalizable the research outcomes are.
% Different types of stations may have different related contexts, issues, and design objectives.
% Nevertheless, the selection of cases in this paper implies a certain station typology in itself. In the scope of this research, the following factors are considered: 

Some features of the selected cases:
1) Continental or national contexts - China and the Netherlands; 
2) Size of stations and cities - the selected cases cover small to large cities and stations (fig. \ref{fig:stationtypes}). Stations' size affects their roles in the cities \citep{vanackerSpatialImpactTrain2021}, and cities' size affects mobility patterns. Size affects the complexity of stations and the cost of spatial-managerial interventions;
3) Security check - which resulted in barriers and potentials for Chinese stations during fluctuations; 
4) Flow pattern - in Chinese stations, the flow of egress and ingress is separated in different paths (like airports), while in Dutch stations, it happens in the same paths; 
5) The way public spaces are used - more diverse leisure activities/events happen in Dutch cases than in Chinese cases.
% (green spaces or not, for CN and NL stations);
% Middle: xxx. Small: Zandvoort, Has a larger fluctuation ratio? Hence insufficient weather-protected area for hot summer? Sizes of station areas: Major stations, middle- and small-size stations (with relevance to fluctuation [goal-oriented], or spatial interventions[means-oriented]).
% - the size of the cities can affect the practicality of the mode. Stations of different sizes can be different in terms of modality (e.g., mode share). Large (Major): Bicycle is an important modality. 
% - the urgency of demand-responsive spatial design. For small stations, management may play a major role without support from spatial design; for major stations, spatial design really matters to reduce managerial costs and interferences between different components. 
% - Layout difference (layout flexibility) due to the variety of elements 
% - this paper only considers function-related issues and leaves aside cultural-related issues.
% - Differences: are or not a monumental cathedral, functionalist machine, or system of infrastructure \citep{vanackerSpatialImpactTrain2021}? - Weaving the tissue of the city?
% It is up to the readers of this paper to decide how far they want to extend the experience gained from the selected cases.


\begin{figure}[h]
    \centering
    \includegraphics[width=0.7\columnwidth]{Paper3Imgs_pdfs/stationtypes.pdf}
    \caption{Station typology implied by the cases}
    \label{fig:stationtypes}
\end{figure} \FloatBarrier







\subsection{Limitations and future research directions}
\label{Limitations and future research directions}

This paper exhibits several limitations and future research directions: 1) There can be hidden limitations of some design principles that we do not yet know. 2) With the same selected cases, there are more pairs of comparisons between them, yet not examined by us due to the project's time frame (We documented to what extent we have examined in \href{https://github.com/1309928130/Paper3_SupplementaryMaterials/blob/main/SupplementaryMaterials.pdf}{Supplementary Materials}, fig. \ref{fig:Examine_relevance_dataprocess}). 3) More cases of different continental and national contexts can be researched in the future. 
% Design vocabulary is also meaningful yet not done.
% (Relates to station typology)

During our interviews with field experts, they pointed out more limitations and future directions: 1) the design principles in this paper are mostly about the spatial form of districts or buildings, while other things are less investigated: programming, user types, spatial qualities, nature-based solutions, and cost-benefit analysis of solutions.







\section{Conclusion}
\label{Conclusions}
In the introduction, we proposed two research questions: 1) What are the design principles for demand-responsive railway station areas? and 2) What is a practical way to develop design principles? 

The first question is answered by the resulting design principles. These design principles are interconnected, covering a wide range of components of the station areas, and can be viewed from various perspectives.

The second question is answered by the seven-step methodology framework for developing design principles. Design principles can be effectively identified through research by design based on multiple data sources. 




\section*{Supplementary Materials}
\label{Supplementary materials}
1. \href{https://github.com/1309928130/Paper3_SupplementaryMaterials/blob/main/SupplementaryMaterials.pdf}{The detailed information} (of the literature, cases, research by design, and evaluations).

2. \href{https://github.com/1309928130/Paper3_SupplementaryMaterials/blob/main/Supplementary%20Material%202%20-%20Sources%20of%20Figures.md}{Sources of figures}.
% \noindent

3. The \href{http://c1309928130.pythonanywhere.com/}{interactive website} for displaying the design principles from different perspectives (\url{http://c1309928130.pythonanywhere.com/})




\setstretch{1} % Adjust this value to your desired line spacing
\setlength{\bibsep}{3pt} % Adjust this value to your desired space between references

\bibliographystyle{elsarticle-harv}
\bibliography{My_Library8}














\setcounter{section}{1}
\newpage
\section*{Supplementary Material 1 - The detailed information}





\subsection{Details of the initial paper searching and selection}

\subsubsection{Enrich the meaning of ‘demand-responsive station areas’ and establish the first theme for literature searching}

Since overcrowding and emptiness are very particular issues that are barely addressed by any existing spatial design research, therefore if directly using 'railway station area’ + ‘demand-responsive’ as the keywords for searching, hardly any relevant papers can be found, and no heuristics can be gained. An enriched meaning of ‘demand-responsive station areas’ is needed to serve as a basis for literature selection. This paper enriches the term’s meaning by reinterpreting the identities of 'railway station areas' and the nature of 'overcrowding' and ‘emptiness.’

The identities of railway stations: According to \citet{bertoliniCitiesRailsRedevelopment1998}, ‘railway stations' are `nodes' of the transport networks and `places' of the cities, so stations are fundamentally `transport infrastructures' for transportation activities and `public spaces' for other leisure activities. 

`Station areas' include the 'station buildings' and the outdoor `spaces,' therefore, ‘open/urban/public space' and ‘landscape' are relevant. `Overcrowding’ can be seen as a type of `risk' to the station system, which can be addressed by `resilience' as a quality of the system. ‘Emptiness’ is a lack of use of public spaces, therefore `flexible' use can be helpful. Both ‘resilience' and `flexible' mean a system that is capable of changing, and similar words can include `dynamic,' `adaptive,' ‘robust,’ `moveable,' `reconfigurable,' `multi-functional (in the time dimension),' `temporary,' `kinetic', and so on. The combination of these concepts facilitates this paper in finding relevant literature (See the main text, Section \ref{Literature} \hyperref[Literature]{Literature}, the first theme of searching) and cases.





\subsubsection{Paper selection criteria}
The papers that used as the heuristics for research-by-design meet at least one of the following criteria: 1) The papers are related to design aspects, including design qualities, programming/functions, layouts, spaces, elements, and so on.
2) The papers deal with demands and users in station areas. 
3) The papers' knowledge is preferably applicable on both building and outdoor spaces, otherwise too specific and not generalizable. In paper selection, we excluded plenty of papers that are purely about planning, or managerial solutions, which can hardly be transformed to spatial interventions.

% Too complicated for design practice?)
% -expert scoring?

% (as flexibility)
% (Exclude the planning, managerial solutions)
% (Ardeshiri et al., 2016) (urban spaces)
% prerequisites: permeability, adaptability, legibility 
% followed by body requisites: positive outdoor spaces, spaces with multimodal behavioral patterns, multiuse buildings, active frontages, and fine amenities 
% constructive elements: activity, place, and people
% Classes: versatility, convertibility, and expansibility

% (Geraedts, 2008) (buildings) (should only be included when meaningful transformations are found)
% Partitionability, adaptability, extendibility, and multifunctionality;
% 1 Partitionable: Partitionable, collective/individual, disconnectable, zonable, modular
% 2 Adaptable: Adaptable, dismantlable, rearrangeable, adjustable, exchangeable, alterable, mobile, shapable
% 3 Extendible: Extendible, capacity, dimensions, ductless
% 4 Multifunctional: Multifunctional, intelligent, automated, universal, integrated
% modularity
% appropriation

% (as resilience)
% (Exclude the managerial qualities)


\subsubsection{Initial papers selected through each searching strategy and their relevance to the layouts and elements of design principles}

\begin{figure}[!ht]
    \centering
    \includegraphics[width=0.6\linewidth]{Paper3Imgs_pdfs/sankeyLR.pdf}
    \caption{The literature selected and the relevance to the research findings}
    \label{fig:sankeyLR}
\end{figure} \FloatBarrier






\newpage
\subsection{The validation of the research by design method}

To some readers who deal with more science-related and quantitative research, this research by design (RbD) method may raise the concern of bias, but in this paper we argue that it is not an issue. This research is for developing mean-end knowledge for designers; As design-related research, it is not about true or false (and therefore the main quality criteria is not heavily rely on validity, reliability, and generalizability, as in quantitative research \citep{leungValidityReliabilityGeneralizability2015}), but about useful/practical or not \citep{zwartMethodologicalClassificationInnovative2016, nijhuisDesignrelatedResearchLandscape2012}. 

Design iterations seem arbitrary as it can iterates numerous times, while the minimum requirement is to iterate until requirements (design objectives) are met. RbD is useful as long as one effective solution is successfully found. 

To further make the data processing clear and transparent, we documented the case comparisons we made in detail, so that readers can check to what extent have we explored with the selected cases, and how much still remain unexplored.

We reflected our outcomes internally by checking all the criteria with each principle, and we did expert interviews to further externally reflect on principles' limitations.


% Design vocabulary
% Design vocabulary is more ambiguous, while design principles are more concrete. Both of them have values in design practice. Paradigms and

% 3.3 Design vocabulary for designing ‘demand-responsive’ station areas (500)
% (Inductive reasoning) (network analysis, Co-occurrence analysis) (meanings of concepts are always evolving) (De Jonge \& van der Voordt, 2002) (Liu, 2020)

% Design vocabulary is useful as an input for generating design principles; it is a byproduct of this research when seeking final design principles.

\vspace{20pt} 

% \newpage
\subsection{The diagrams of station areas}

As shown in the following diagrams, and the on-site visits, the station areas are hardly perceivable (by vision) outside the 250m range.


\begin{figure}[!ht]
    \centering
    \includegraphics[width=\linewidth]{Paper3Imgs_pdfs/2024Diagrams.001.pdf}
    \caption{The Utrecht Central Station area - 250m vs 500m}
    \label{fig:station_area_Utrecht}
\end{figure} \FloatBarrier

\begin{figure}[!ht]
    \centering
    \includegraphics[width=\linewidth]{Paper3Imgs_pdfs/2024Diagrams.003.pdf}
    \caption{The Rotterdam Central Station area - 250m vs 500m}
    \label{fig:station_area_Rotterdam}
\end{figure} \FloatBarrier

\begin{figure}[!ht]
    \centering
    \includegraphics[width=\linewidth]{Paper3Imgs_pdfs/2024Diagrams.002.pdf}
    \caption{The Beijing West Railway Station area - 250m vs 500m}
    \label{fig:station_area_BeijingWest}
\end{figure} \FloatBarrier



\clearpage
\newpage

\subsection{Initial station cases selected}

\begin{table*}[ht]
\centering
\caption{Cases selected}
\label{tab:case_selected}
    \begin{tabular}{p{4.5cm}p{1.5cm}p{2.5cm}p{2.5cm}p{3cm}}
\toprule
    \textbf{Stations} & \textbf{Size} & \textbf{Demand fluctuation} & \textbf{Performance for fluctuation} & \textbf{Nation} \\
\midrule
    Beijing West RS &	Large &	Large &	Poor&	China \\
    Beijing South RS&	Large& 	Large &	Poor&	China \\
    Shanghai RS& 	Large& 	Large& 	Good &	China\\
    Rotterdam Central Station& 	Large &	Middle	&Good	&The Netherlands\\
    Utrecht Central Station &	Large &	Middle	&Good&	The Netherlands\\
    Bijlmer Arena RS &	Middle	&Large &	Good&	The Netherlands\\
    Zandvoort Station &	Small	&Large &	Good&	The Netherlands\\
\bottomrule
    \end{tabular}
\end{table*}









\clearpage
\subsection{Data processing}


% \rotatebox[origin][width][height]{angle}{content}
\begin{figure}[!htbp]
    \centering
    \rotatebox{90}{\includegraphics[width=0.9\textheight]{Paper3Imgs_pdfs/dataprocess.pdf}}
    \caption{Examine relevance and compare cases}
    \label{fig:Examine_relevance_dataprocess}
\end{figure} \FloatBarrier

% Notes: Some comparisons are not examined, as shown in the following table. 


\clearpage
\subsection{The table for internal evaluation of design principles}
% \clearpage

\begin{figure}[!ht]
    \centering
    \includegraphics[width=0.9\linewidth]{Paper3Imgs_pdfs/evaluateDP.pdf}
    \caption{The evaluation of design principles}
    \label{fig:evaluateDP}
\end{figure} \FloatBarrier




\clearpage
\subsection{Relate design principles to an event typology}

\textcolor{red}{(Smith et al., 2021) based on access and mobility (relevance: for spatial-managerial consideration and discussion); in station areas or not? Event relevance: -(access+mobility) Interactions/conflicts with the (passenger and non-passenger) flow or not -Size. relationship between paths and event types.
No significant influence on overcrowding. Layout matters in the sense of: xxxx, . As long as it makes sense to designers who use these solutions.}

\begin{figure}[!htbp]
    \centering
    \includegraphics[width=0.75\linewidth]{Paper3Imgs_pdfs/event_principles.pdf}
    \caption{The design principles related to the event typology}
    \label{fig:eventtypology_principles}
\end{figure} \FloatBarrier





\clearpage
\newpage
\section*{Supplementary Material 2 - Sources of figures}

% \listoffigures

% \begin{xltabular}{\textwidth}{L L L}
\begin{xltabular}{\textwidth}{p{0.01\textwidth} p{0.42\textwidth} p{0.50\textwidth}} % adjust widths manually

    \caption{Sources of Figures} 
    \label{tab:figsources} \\
    
    \hline
    \textbf{N} & \textbf{Figure and Caption} & \textbf{Source} \\
    \hline
    \endfirsthead
    
    \caption{Sources of Figures (continued)} \\
    \hline
    \textbf{N} & \textbf{Figure and Caption} & \textbf{Source} \\
    \hline
    \endhead
    
% Table content starts here
\ref{fig:overcrowdingemptiness} & Overcrowding and emptiness in railway stations areas. &  Left - \url{https://www.163.com/dy/article/DLED9I020524TFU4.html}; Right - \url{https://www.sohu.com/a/386779398_640391}\\

\ref{fig:flowunit} & A conceptual basic unit of passenger flow. & By authors \\

\ref{fig:bottleneck} & The mechanism of overcrowding and emptiness in Beijing West railway station area. & By authors \\

\ref{fig:beijingwestboth} & Spatial relevance shown by an example case of Beijing West Railway Station. & Original photo sources: Left - \url{https://www.meipian.cn/}; Right - \url{https://upload.wikimedia.org/wikipedia/commons/5/5b/Beijing_West_Railway_Station_20170506_124557.jpg}. Drawings by authors \\

\ref{fig:temporaldimension} & Temporal scales. & By authors, inspired by \citep[p. 38, fig. 13]{dejongeWAYSStudyResearch2002} \\

% \ref{fig:Spatial scale of events} & Spatial scales of events. & By authors \\

\ref{fig:spatial_scales_station} & Spatial scales of station areas. & By authors \\

\ref{fig:Archi_urban_proposals_Sloterdijk} & Architectural and urban design proposals for Amsterdam Sloterdijk station. & Upper - \citep{triggianeseStationsNodesExploring2019}; Lower - \citep{andrianosNodePlaceNew2023}\\

\ref{fig:positionP3inPhD} & This paper is part of a PhD project at the intersection of three topics. & By authors \\

\ref{fig:methodologyframework} &  The process of developing design principles. & By authors \\

\ref{fig:configuration} &  The station at the district and building levels. & By authors \\

\ref{fig:configurationflow} & The configuration of the station. & By authors \\

\ref{fig:datacompare} & Data sources. & By authors \\

\ref{fig:researchbydesign} & An example of research by design. & By authors \\

\ref{fig:sourcesofheuristics} & Sources of heuristics. & By authors \\

\ref{fig:Pattern_language_network} & The design principles (patterns) within a network. & By authors \\

\ref{fig:perspectivematrix} & The design principles viewed from different perspectives. & By authors \\



\ref{fig:flexible_use} & Flexible use. & \url{https://www.stadsstranden.nl/paris-plages/} \\

\ref{fig:generallayout} & Different general layouts of the station and city. & \citep{qiAnalysisStationCitySynergetic2019} \\

\ref{fig:utrechtsetapartnonstation} & Set apart non-transport function. & \url{https://www.santenco.nl/portfolio_page/stationsplein-oost/} \\

\ref{fig:districtscattered} & Vibrate city environment by scattered mobility nodes. & By authors \\

\ref{fig:buildingscattered} & Increase transport capacity by scattered mobility nodes. & By authors \\

\ref{fig:pathwithevent} & Alignment between open spaces and main paths. & By authors \\

\ref{fig:humancenteredspace} & Human-oriented spaces versus vehicle-oriented spaces. & By authors \\

\ref{fig:Rot_old_new} & The old and the new Rotterdam central stations. & Upper-left - \url{https://commons.wikimedia.org/wiki/File:Nieuw_Centraal_Station_in_Rotterdam_aangetast,_Bestanddeelnr_908-5455.jpg}; Upper-right -\url{https://internaathetposthuis.nl/raldus.html};  Lower-left -\url{https://indebuurt.nl/rotterdam/nieuws/salsa-doe-gratis-mee-met-deze-gigadans-voor-rotterdam-centraal~311821/}; Lower-right \url{https://mvsa-architects.com/en/news/nos-rotterdam-central-receives-a-score-of-8-and-is-the-most-popular-large-train-station/}\\

\ref{fig:connection} & Connect with neighborhoods. & By authors \\

\ref{fig:jiaxing} & Smooth level changes by landscape design at the district level. & \url{https://m.thepaper.cn/newsDetail_forward_5348462} \\

\ref{fig:arnhem} & Smooth level changes by landscape design at the building level. & \url{https://www.unstudio.com/en/page/12109/arnhem-central-masterplan} \\

\ref{fig:fencezigzag} & Adaptive redundant spaces with path regulation. & By authors \\

\ref{fig:Path_regula_reconfig_elements} & Path regulation using reconfigurable elements. & By authors \\

\ref{fig:programmingsharing} & Programming considering the time dimension (for flexible use). & By authors \\

\ref{fig:buildshortcut} & Shortcuts or optimizing paths. & By authors \\

\ref{fig:theshed} & A changeable building - the Shed. & \url{https://www.architecturaldigest.com/story/the-shed-finally-opens-new-york-city-hudson-yards} \\

\ref{fig:Changeable_building_components} & Changeable building components. & By authors \\

\ref{fig:antwerpferrywheel} & Add installations and facilities. & \url{https://www.flickr.com/photos/87453322@N00/7561351848} \\

\ref{fig:urbanfurniture} & Reconfigurable elements. & \citep{iversStagingUrbanLandscapes2018} \\

\ref{fig:Reconfigurable_spaces} & Reconfigurable spaces. & Left - \url{https://www.shutterstock.com/image-photo/view-delft-market-square-nieuwe-kerk-1206141775}; Right - \url{https://nl.pinterest.com/pin/546342998538754823/} \\

\ref{fig:layout_setapart} & Redundant spaces or setting apart bottlenecks at the district level. & By authors \\

\ref{fig:build_setapart} & Redundant spaces or setting apart bottlenecks at the building level. & By authors \\

\ref{fig:layoutpassage} & Position city passages (at the district level). & By authors \\

\ref{fig:buildpassage} & The city passages (at the building level). & \url{https://www.amsterdam-viptours.com/blog/cuypers-passage/} \\

\ref{fig:beijingwestnorth} & Barriers that hinder the flow at the district level. & By authors \\

\ref{fig:buildflowcommercial} & Reduce barriers to ease flow at the building level. & By authors \\

\ref{fig:utrechtstair} & Stairs as stages or seats. & \url{https://architectenweb.nl/nieuws/artikel.aspx?id=43948} \\

\ref{fig:City axes} & Two city axes were planned during the redevelopment of the Utrecht Central Station area. & \href{https://www.bouwpututrecht.nl/wp-content/uploads/2015/12/12A-OO-Structuurplan-Stationsgebied-2006.pdf}{Structuurplan, Stationsgebied Utrecht, December 2006} \\

\ref{fig:stationtypes} & Station typology implied by the cases. & By authors \\








\ref{fig:station_area_Utrecht} & The Utrecht Central Station area - 250m vs 500m. & By authors \\

\ref{fig:station_area_Rotterdam} & The Rotterdam Central Station area - 250m vs 500m \\

\ref{fig:station_area_BeijingWest} & The Beijing West Railway Station area - 250m vs 500m. & By authors \\

\ref{fig:sankeyLR} & The literature selected and the relevance to the research findings. & By authors \\


\ref{fig:Examine_relevance_dataprocess} & Examine relevance and compare cases. & By authors \\

\ref{fig:evaluateDP} & The evaluation of design principles. & By authors \\

\ref{fig:eventtypology_principles} & The design principles related to the event typology. & By authors \\



    \hline
\end{xltabular}

\end{document}



\noindent
\textbf{Note}: All other figures without mentioning sources are drawn by the authors of this paper.





\end{document}
